\exercisesection{}

\begin{enumerate}
\def\labelenumi{\arabic{enumi})}
\item
  Let G be a locally compact group, H a closed subgroup of G. For \(\xi\in H\) , set \(\chi(\xi)=\Delta_{H}(\xi)\Delta_{G}(\xi)^{-1}\) . Regard H as operating on the right in G by translations. Show that there exists on \(\mathcal{K}^{\chi}(G)\) a nonzero positive linear form I and, up to a constant factor, only one, that is invariant under left translation. (Make use of Prop. 3 with X=G, while taking \(\mu\) to be a left Haar measure of G.)
\item
  Let X and \(X'\) be two locally compact spaces in which a locally compact group H operates on the right, continuously and properly. Let \(\theta\) be a proper continuous mapping of X into \(X'\) , compatible with the identity mapping of H (GT, III, §2, No.~4), and let \(\theta': X/H \to X'/H\) be the mapping deduced from \(\theta\) by passage to quotients. Let \(f'\) be a continuous function on \(X'\) whose support has compact intersection with the saturation of each compact set. Then \(f = f' \circ \theta\) has the same properties in X, and
\end{enumerate}

\[
f ^ {b} = f ^ {\prime b} = \theta^ {\prime}
\]

(the mappings \(f \mapsto f^b\) , \(f' \mapsto f'^b\) being relative to the choice of a same Haar measure on H).

\begin{enumerate}
\def\labelenumi{\arabic{enumi})}
\setcounter{enumi}{2}
\tightlist
\item
  Let B be a locally compact space, H a locally compact group. Set \(\mathbf{X} = \mathbf{B} \times \mathbf{H}\) , the group H operating in X by
\end{enumerate}

\[
(b, \xi) \xi^ {\prime} = (b, \xi \xi^ {\prime}).
\]

Let \(\lambda\) be a measure on \(B = X / H\) , and \(\beta\) a left Haar measure on \(H\) . Then \(\lambda^{\sharp} = \lambda \otimes \beta\) .

\begin{enumerate}
\def\labelenumi{\arabic{enumi})}
\setcounter{enumi}{3}
\item
  Let X be a locally compact space in which a locally compact group H acts on the right continuously and properly. Let \(\beta\) be a left Haar measure on H, \(\mu\) a measure \(\geqslant0\) on X. Let h be a continuous function \(\geqslant0\) on X such that \(h^{b}=1\) . If f is a \(\mu\) -negligible function \(\geqslant0\) on X, the function \((x,\xi)\mapsto f(x)h(x\xi)\) is \((\mu\otimes\beta)\) -negligible on X×H. (There exist decreasing open sets \(\Omega_{1},\Omega_{2},\ldots\) such that \(\mu(\Omega_{n})\) tends to 0 and such that f is zero outside \(\Omega_{1}\cap\Omega_{2}\cap\cdots\) . Show that \(\int^{*}\varphi_{\Omega_{n}}(x)h(x\xi)d\mu(x)d\beta(\xi)\leqslant\mu(\Omega_{n})\) , on observing that the function \((x,\xi)\mapsto\varphi_{\Omega_{n}}(x)h(x\xi)\) is lower semi-continuous.)
\item
  Let \(\mathbf{G}\) be a locally compact group. For every \(s \in \mathbf{G}\) , let \(\psi_s\) be the automorphism of the additive group \(\mathbf{R}\) defined by \(\psi_s(x) = \Delta_{\mathbf{G}}(s)x\) . Let \(\Gamma\) be the topological semi-direct product of \(\mathbf{G}\) and \(\mathbf{R}\) defined by \(s \mapsto \psi_s\) (GT, III, §2, No.~10). Show that \(\Gamma\) is unimodular.
\item
  Let \(G\) be a locally compact group, \(G_1, G_2, G_3\) closed subgroups such that \(G_3 \subset G_1 \cap G_2\) . Assume that \(G, G_1, G_2, G_3\) are unimodular, and that \(G_1 / G_3\) has finite total measure (for every measure invariant under \(G_1\) ). Let \(\lambda, \mu, \nu\) be invariant measures on \(G / G_1, G / G_2, G_2 / G_3\) , and \(\varphi\) the canonical mapping of \(G\) onto \(G / G_1\) . Let \(f \in \mathcal{K}(G / G_1)\) . Show that
\end{enumerate}

\[
a \int_ {\mathrm{G} / \mathrm{G} _ {1}} f (u) d \lambda (u) = \int_ {\mathrm{G} / \mathrm{G} _ {2}} d \mu (\dot {x}) \int_ {\mathrm{G} _ {2} / \mathrm{G} _ {3}} f (\varphi (x \xi)) d \nu (\dot {\xi})
\]

\((\dot{x} = x\mathrm{G}_{2}, \dot{\xi} = \xi \mathrm{G}_{3})\) , where \(a\) is a constant independent of \(f\) .

¶ 7) a) Let E be a locally compact space, Γ a locally compact group operating on the left continuously in E. Assume that for every \(x \in E\) , the mapping \(s \mapsto s \cdot x\) of Γ into E is proper, and that there exists a nonzero bounded positive measure μ on E invariant under Γ. Then Γ is compact. (Let \((s_n)\) be a sequence of points of Γ. Let K be a compact subset of E that is not μ-negligible. Show, using Exer. 15 of Ch. IV, §4, that there exists an \(x_0 \in E\) such that \(s_n x_0 \in K\) for infinitely many values of n.~From this, deduce that \((s_n)\) has a cluster point in Γ. Then apply Exer. 6 of GT, II, §4.)

\begin{enumerate}
\def\labelenumi{\alph{enumi})}
\setcounter{enumi}{1}
\tightlist
\item
  Let G be a locally compact group, H and K two closed subgroups of G. Assume that G and H are unimodular and that G/H has finite measure for the measure invariant under G. Show that for H and K to satisfy the equivalent conditions of Exer. 11 c) of GT, III, §4, it is necessary and sufficient that K be compact.
\end{enumerate}

\begin{enumerate}
\def\labelenumi{\arabic{enumi})}
\setcounter{enumi}{7}
\tightlist
\item
  Let \(\mathfrak{G}\) be a locally compact group, G and H two closed subgroups of \(\mathfrak{G}\) . Assume that every \(s \in \mathfrak{G}\) can be written in one and only one way in the form \(s = \xi x = y\eta\) \((\xi, \eta \in \mathbf{G}, x, y \in \mathbf{H})\) , and that \(\xi, x, y, \eta\) depend continuously on \(s\) . Every \(\xi \in \mathbf{G}\) defines a homeomorphism \(\widehat{\xi}\) of H onto H by \(\xi x \in \widehat{\xi}(x)\mathbf{G} (x \in \mathbf{H})\) . Every \(x \in \mathbf{H}\) defines a homeomorphism \(\widehat{x}\) of G onto G by \(\xi x \in \mathrm{H}\widehat{x}(\xi) (\xi \in \mathbf{G})\) . Let \(\mu, \alpha, \beta\) be left Haar measures on \(\mathfrak{G}, \mathbf{G}, \mathbf{H}\) . Show that
\end{enumerate}

\[
\begin{array}{c} d \widehat {x} ^ {- 1} (\alpha) (\xi) = \frac {\Delta_ {\mathfrak {G}} (\widehat {\xi} (x)) \Delta_ {\mathrm{G}} (\widehat {x} (\xi))}{\Delta_ {\mathrm{H}} (\widehat {\xi} (x)) \Delta_ {\mathrm{G}} (\xi)} d \alpha (\xi), \\ d \widehat {\xi} ^ {- 1} (\beta) (x) = \frac {\Delta_ {\mathrm{G}} (\widehat {x} (\xi))}{\Delta_ {\mathfrak {G}} (\widehat {x} (\xi))} d \beta (x). \end{array}
\]

(Let \(f \in \mathcal{K}(G)\) . Express the fact that \(\int f(u) d\mu(u)\) , calculated with the help of (23), either for \(X = G\) , \(Y = H\) , or for \(X = H\) , \(Y = G\) , is invariant when \(f\) is subjected to left translation by an element of \(G\) or an element of \(H\) .)

\begin{enumerate}
\def\labelenumi{\arabic{enumi})}
\setcounter{enumi}{8}
\tightlist
\item
  Let X be a locally compact space, H a compact group operating continuously (hence properly) on the right in X. Denote by \(\beta\) the normalized Haar measure on H, by \(\pi: X \to X/H\) the canonical mapping. Let \(\lambda\) be a positive measure on X/H, \(\lambda^{\sharp}\) the corresponding measure on X.
\end{enumerate}

\begin{enumerate}
\def\labelenumi{\alph{enumi})}
\item
  Show that if N is a \(\lambda\) -negligible subset of X/H, then \(\overline{\pi}^{1}(N)\) is \(\lambda^{\sharp}\) -negligible. (Calculate the measure of a set \(\overline{\pi}^{1}(U)\) , where U is open in X/H.)
\item
  Let \(p\) be a finite number \(\geqslant 1\) , and let \(f\) be a function in \(\mathcal{L}_{\mathrm{F}}^{p}(\mathrm{X},\lambda^{\sharp})\) (F a Banach space or \(\mathrm{F} = \overline{\mathbf{R}}\) ). Show that the set of \(x \in \mathrm{X}\) such that \(\xi \mapsto f(x\xi)\) is not \(\beta\) -integrable on \(\mathrm{H}\) is of the form \(\overline{\pi}^{1}(\mathrm{N})\) , where \(\mathrm{N}\) is \(\lambda\) -negligible. Moreover, if \(f^b\) is the function on \(\mathrm{X / H}\) , defined almost everywhere (for \(\lambda\) ), such that \(f^b (\pi (x)) = \int_{\mathrm{H}} f(x\xi) d\beta (\xi)\) , then \(f^b\) belongs to \(\mathcal{L}_{\mathrm{F}}^{p}(\mathrm{X / H},\lambda)\) , and \(\mathrm{N}_p(f^b)\leqslant \mathrm{N}_p(f)\) .
\item
  Conversely, if \(g \in \mathcal{L}_{\mathrm{F}}^{p}(\mathrm{X / H}, \lambda)\) , then the function \(g \circ \pi\) belongs to \(\mathcal{L}_{\mathrm{F}}^{p}(\mathrm{X}, \lambda^{\sharp})\) . If \(p, q\) are conjugate exponents, \(\mathrm{F}'\) is the dual of \(\mathrm{F}\) , \(f \in \mathcal{L}_{\mathrm{F}}^{p}(\mathrm{X}, \lambda^{\sharp})\) and \(g \in \mathcal{L}_{\mathrm{F}'}^{q}(\mathrm{X / H}, \lambda)\) , then
\end{enumerate}

\[
\int_ {\mathrm{X}} \left\langle f (x), g (\pi (x)) \right\rangle d \lambda^ {\sharp} (x) = \int_ {\mathrm{X/H}} \left\langle f ^ {\flat} (z), g (z) \right\rangle d \lambda (z).
\]

\begin{enumerate}
\def\labelenumi{\arabic{enumi})}
\setcounter{enumi}{9}
\tightlist
\item
  With the notations of §1, No.~6, Prop. 6 let, for every \(\alpha \in A\) , \(\mu_{\alpha}\) be a left Haar measure on \(G_{\alpha}\) , so that the inverse limit of the \(\mu_{\alpha}\) is a left Haar measure \(\mu\) on \(G\) . For every \(\alpha \in A\) , denote by \(\lambda_{\alpha}\) the normalized Haar measure on \(K_{\alpha}\) .
\end{enumerate}

\begin{enumerate}
\def\labelenumi{\alph{enumi})}
\tightlist
\item
  Let \(p\) be a finite number \(\geqslant 1\) and let \(f\) be a function in \(\mathcal{L}_{\mathrm{F}}^{p}(\mathrm{G},\mu)\) (F a Banach space or \(\mathrm{F} = \overline{\mathbf{R}}\) ). For every \(\alpha \in \mathbf{A}\) , the function
\end{enumerate}

\[
f _ {\alpha} (s) = \int_ {K _ {\alpha}} f (s \xi) d \lambda_ {\alpha} (\xi),
\]

defined almost everywhere for \(\mu\) , belongs to \(\mathcal{L}_{\mathrm{F}}^{p}(\mathrm{G},\mu)\) (Exer. 9). Show that, with respect to the directed set A, \(f_{\alpha}\) tends in mean of order \(p\) to \(f\) (make use of Lemma 2 of §1, No.~6).

\begin{enumerate}
\def\labelenumi{\alph{enumi})}
\setcounter{enumi}{1}
\tightlist
\item
  Assume that A = N and p = 1. Show that \(f_{n}\) then tends to f almost everywhere (for \(\mu\) ) in G (use a method analogous to that of Ch. V, §8, Exer. 16).
\end{enumerate}

¶ 11) Let G be a locally compact group, H a closed subgroup of G, λ a left Haar measure on G; assume that there exists on G/H a measure μ invariant under G such that λ = μ \(^{\sharp}\) and μ(G/H) \textless{} +∞. Let ν be a nonzero positive measure on G/H such that \(\nu (\mathrm{G / H}) < +\infty\) . Finally, let \(h\) be a function on \(\mathbf{G}\) having the properties of No.~4, Prop. 8.

\begin{enumerate}
\def\labelenumi{\alph{enumi})}
\tightlist
\item
  Let A be a Borel subset of G/H. Show that
\end{enumerate}

\[
\int_ {\mathrm{G}} h (s) \nu (s ^ {- 1} \mathrm{A}) d \lambda (s) = \nu (\mathrm{G/H}) \mu (\mathrm{A})
\]

(make use of Prop. 9 a) of No.~4).

\begin{enumerate}
\def\labelenumi{\alph{enumi})}
\setcounter{enumi}{1}
\tightlist
\item
  Let \(A_i\) ( \(1 \leqslant i \leqslant n\) ) be Borel subsets of \(G / H\) and, for every \(i\) , let \(b_i = \sup_{s \in G} \nu(s^{-1} A_i)\) . Show that there exist two distinct indices \(i, j\) such that
\end{enumerate}

\[
\int_ {\mathrm{G}} h (s) \nu (s ^ {- 1} \mathrm{A} _ {i}) \nu (s ^ {- 1} \mathrm{A} _ {j}) d \lambda (s) \geqslant \frac {\left(\nu (\mathrm{G} / \mathrm{H})\right) ^ {2}}{\mu (\mathrm{G} / \mathrm{H})} \left(\left(\sum_ {k = 1} ^ {n} \mu (\mathrm{A} _ {k})\right) ^ {2} - \frac {\mu (\mathrm{G} / \mathrm{H})}{\nu (\mathrm{G} / \mathrm{H})} \sum_ {k = 1} ^ {n} b _ {k} \mu (\mathrm{A} _ {k})\right)
\]

(argue as in §1, Exer. 28 a)).

¶ 12) a) Let X be a topological space, \(f: X \to N\) an upper semi-continuous function. Let \(X_0\) be the set of points of X where \(f\) is locally constant, and \(Y_0 = X - X_0\) . Define recursively two sequences \((X_i)_{i \geqslant 0}\) , \((Y_i)_{i \geqslant 0}\) of subsets of X, in the following way: \(X_i\) is the set of points of \(Y_{i-1}\) where \(f|_{Y_{i-1}}\) is locally constant, and \(Y_i = Y_{i-1} - X_i\) . Show that \(X_i\) is a dense open subset of \(Y_{i-1}\) , and that \(\bigcap_i Y_i = \varnothing\) .

(Show that \(f(x) > i\) at every point of \(Y_i\) .)

\begin{enumerate}
\def\labelenumi{\alph{enumi})}
\setcounter{enumi}{1}
\tightlist
\item
  Let X be a set, H a group operating on the right in X, \(\pi\) the canonical mapping of X onto X/H, and A (resp. B) the set of A ⊂ X such that \(\pi|A:A\to X/H\) is injective (resp. surjective). Let \((V_{1},V_{2},\ldots)\) be a countable covering of X by elements of A. Let
\end{enumerate}

\[
\mathrm{V} _ {i} ^ {\prime} = \mathrm{V} _ {i} \cap \complement \left(\mathrm{V} _ {1} \mathrm{H} \cup \mathrm{V} _ {2} \mathrm{H} \cup \dots \cup \mathrm{V} _ {i - 1} \mathrm{H}\right).
\]

Show that \(\mathrm{F} = \bigcup_{i\geqslant 0}\mathrm{V}_i^\prime \in \mathfrak{A}\cap \mathfrak{B}\) .

\begin{enumerate}
\def\labelenumi{\alph{enumi})}
\setcounter{enumi}{2}
\item
  Let X be a locally compact space, H a countable discrete group operating on the right in X continuously and properly, \(\pi\) the canonical mapping of X onto X/H, and \(\mu\) a measure \(\geqslant 0\) on X. Show that if the sets \(V_{i}\) of b) are quadrable (Ch. IV, §5, Exer. 17 d)), then the set F of b) is a quadrable fundamental domain.
\item
  Let us maintain the hypotheses of \(c\) ), and use the notations \(\mathbf{H}_x, n(x)\) of No.~10. An \(x \in X\) is said to be general if there exists a neighborhood \(V\) of \(x\) in \(X\) such that \(\mathrm{H}_y = \mathrm{H}_x\) for all \(y \in V\) . Show that the general points of \(X\) are those where the function \(n\) is locally constant. Show that a general point admits an open neighborhood that belongs to \(\mathfrak{A}\) and is quadrable (make use of Ch. IV, §5, Exer. 17 d)).
\item
  Let us maintain the hypotheses of \(c\) ), and assume moreover that \(X\) is countable at infinity. Show that if \(X - X_0\) is \(\mu\) -negligible, then there exists a fundamental domain \(F\) that is Borel and quadrable. (Apply the construction of \(a\) ) to the function \(n\) . Then apply the results of \(c\) ) and \(d\) ) in \(X_0\) . Argue similarly in each \(X_i\) .)
\item
  Let \(U \in B\) be an open set. Show that one can impose on the set F of e) the following supplementary properties: 1) \(F \subset U\) ; 2) for every compact subset K of X, the set of \(s \in H\) such that Fs intersects K is finite.
\end{enumerate}

\begin{enumerate}
\def\labelenumi{\arabic{enumi})}
\setcounter{enumi}{12}
\tightlist
\item
  Let G be a compact group, \(\mu\) a Haar measure on G, u an endomorphism of G such that \(u(\mathrm{G})\) is an open subgroup of G and the kernel \(\overline{u}(e)\) (denoted \(G_{u}\) ) a finite subgroup of G.
\end{enumerate}

\begin{enumerate}
\def\labelenumi{\alph{enumi})}
\item
  Show that there exist a real number \(h(u) > 0\) and an open neighborhood \(U\) of \(e\) in \(G\) such that, for every open set \(V \subset U\) , \(u(V)\) is open in \(G\) and \(\mu(u(V)) = h(u)\mu(V)\) (cf.~§1, No.~7, Cor. of Prop. 9).
\item
  Show that \(h(u) = \operatorname{Card}(\mathrm{G} / u(\mathrm{G})) / \operatorname{Card}(\mathrm{G}_u)\) (calculate \(\mu(u(\mathrm{G}))\) in two ways, using a) and Prop. 10 of No.~7).
\end{enumerate}

\exercisesection{}

\begin{enumerate}
\def\labelenumi{\arabic{enumi})}
\item
  Let E be a finite-dimensional vector space over R, C or H, \(\Phi\) and \(\Phi'\) two nondegenerate positive hermitian forms on E. Assume that \(\mathbf{U}(\Phi) \subset \mathbf{U}(\Phi')\) . Show that \(\Phi\) and \(\Phi'\) are proportional. (Make use of the fact that there exists an orthonormal basis \((e_{1}, \ldots, e_{n})\) for \(\Phi\) that is orthogonal for \(\Phi'\) . The mapping \(u \in \mathcal{L}(\mathrm{E}, \mathrm{E})\) such that \(u(e_{i}) = e_{j}\) , \(u(e_{j}) = e_{i}\) , \(u(e_{k}) = e_{k}\) for \(k \neq i, j\) belongs to \(\mathbf{U}(\Phi)\) .)
\item
  Adopt the notations of Lemma 7. Show that \(\Omega\) has negligible complement in \(\mathbf{GL}(n,\mathbf{K})\) for the Haar measure. (Argue as for Prop. 6.)
\end{enumerate}

¶ 3) Let X be a locally compact space in which a locally compact group H operates on the right, continuously and properly, by \((x,\xi)\mapsto x\xi\) ( \(x\in X\) , \(\xi\in H\) ). Let \(\pi\) be the canonical mapping \(X\to X/H\) . Let \(\rho\) be a continuous representation of H in \(\mathbf{GL}(n,\mathbf{C})\) . Show that for every \(b\in X/H\) , there exist a neighborhood U of b in X/H and a continuous mapping r of \(\overline{\pi}(\mathrm{U})\) into \(\mathbf{GL}(n,\mathbf{C})\) such that \(r(x\xi)=r(x)\rho(\xi)\) for all \(x\in\overline{\pi}(\mathrm{U})\) and \(\xi\in H\) . (One may assume X/H to be compact. Let f be a continuous function \(\geqslant0\) on X, not identically zero on any orbit, and with compact support. Set

\[
r (x) = \int_ {\mathbf {H}} f (x \xi) \rho (\xi) ^ {- 1} d \beta (\xi) \in \mathbf {M} _ {n} (\mathbf {C}),
\]

where \(\beta\) denotes a left Haar measure on H. Show that, for \(f\) and U suitably chosen, \(r(x) \in \mathbf{GL}(n, \mathbf{C})\) for every \(x \in \overline{\pi}^1(\mathrm{U})\) .

\begin{enumerate}
\def\labelenumi{\arabic{enumi})}
\setcounter{enumi}{3}
\tightlist
\item
  Let \(\mathbf{K}\) be a nondiscrete commutative locally compact field. Let \(\mathbf{A}\) be an algebra of finite rank over \(\mathbf{K}\) .
\end{enumerate}

\begin{enumerate}
\def\labelenumi{\alph{enumi})}
\tightlist
\item
  Let \(\mathrm{T}(n, \mathbf{A})\) be the algebra of matrices \((x_{ij}) \in \mathbf{M}_n(\mathbf{A})\) such that \(x_{ij} = 0\) for \(i > j\) . Show that if \(M = (m_{ij}) \in \mathrm{T}(n, \mathbf{A})\) , then
\end{enumerate}

\[
\mathrm{N} _ {\mathrm{T} (n, \mathrm{A}) / \mathrm{K}} (M) = \prod_ {i = 1} ^ {n} \mathrm{N} _ {\mathrm{A} / \mathrm{K}} (m _ {i i} ^ {i - n + 1}).
\]

(Use Lemma 6.)

\begin{enumerate}
\def\labelenumi{\alph{enumi})}
\setcounter{enumi}{1}
\item
  Let \(\mathrm{T}(n,\mathrm{A})^*\) be the group of invertible elements of \(\mathrm{T}(n,\mathrm{A})\) . Let \(M = (m_{ij}) \in \mathrm{T}(n,\mathrm{A})\) . Show that \(M \in \mathrm{T}(n,\mathrm{A})^*\) if and only if \(m_{ii}\) is invertible in \(\mathbf{A}\) for all \(i\) .
\item
  Show that a left Haar measure on \(\mathbf{T}(n,\mathbf{A})^{*}\) is given by
\end{enumerate}

\[
\left(\prod_ {i = 1} ^ {n} \operatorname{mod} \mathrm{N} _ {\mathrm{A} / \mathrm{K}} (m _ {i i}) ^ {i - n - 1}\right) \cdot \bigotimes_ {i \leqslant j} d \alpha (m _ {i j}),
\]

where \(\alpha\) denotes a Haar measure of the additive group of A.

\begin{enumerate}
\def\labelenumi{\alph{enumi})}
\setcounter{enumi}{3}
\tightlist
\item
  Show that if \(\mathbf{N}_{\mathbf{A} / \mathbf{K}} = \mathbf{N}_{\mathbf{A}^0 /\mathbf{K}}\) , then
\end{enumerate}

\[
\Delta_ {\mathrm{T} (n, \mathrm{A}) ^ {*}} \left(\left(m _ {i j}\right)\right) = \prod_ {i = 1} ^ {n} \operatorname{mod} \mathrm{N} _ {\mathrm{A} / \mathrm{K}} \left(m _ {i i}\right) ^ {2 i - n - 1}.
\]

\begin{enumerate}
\def\labelenumi{\arabic{enumi})}
\setcounter{enumi}{4}
\tightlist
\item
  Equip \(\mathbf{R}^{n}\) with the quadratic form \(x_{1}^{2} + x_{2}^{2} + \cdots + x_{n}^{2}\) . Let \(G_{n}\) be the group of displacements of \(\mathbf{R}^{n}\) (Alg., Ch. IX, §6, No.~6, Def. 3) of determinant 1.
\end{enumerate}

\begin{enumerate}
\def\labelenumi{\alph{enumi})}
\item
  Show that \(\mathbf{G}_n\) is unimodular.
\item
  Every element of \(\mathbf{G}_2\) may be uniquely written in the form
\end{enumerate}

\[
(x, y) \mapsto (u + x \cos \omega - y \sin \omega , v + x \sin \omega + y \cos \omega),
\]

where \(u, v\) are in \(\mathbf{R}\) , and \(\omega\) is an element of the group \(\Theta\) of angles of half-lines, isomorphic to \(\mathbf{U}\) . Show that \(du \otimes dv \otimes d\omega\) (where \(d\omega\) is a Haar measure on \(\Theta\) ) is a Haar measure on \(G_2\) .

\begin{enumerate}
\def\labelenumi{\arabic{enumi})}
\setcounter{enumi}{5}
\tightlist
\item
  Let \(\mathbf{P}\) be the set of complex numbers \(z = x + iy\) whose imaginary part is \(>0\) . Show that \(\mathbf{SL}(2,\mathbf{R})\) operates transitively and continuously on the left in \(\mathbf{P}\) by
\end{enumerate}

\[
\left( \begin{array}{c c} a & b \\ c & d \end{array} \right) \cdot z \mapsto \frac {a z + b}{c z + d}
\]

and that P thus becomes a topological homogeneous space for \(\mathbf{SL}(2,\mathbf{R})\) . Show that \(y^{-2}dx dy\) is an invariant measure on P.

¶ 7) Let G be the unimodular group \(\mathbf{SL}(n,\mathbf{R})\) , \(\Gamma\) the subgroup of G formed by the matrices of determinant 1 with integer entries. Show, by induction on n, that every invariant measure \(\mu\) on the homogeneous space \(G/\Gamma\) is bounded and that, for \(f \in \mathcal{K}(\mathbf{R}^{n})\) , one has, with \(\mu\) suitably chosen,

\[
\int_ {\mathbf {R} ^ {n}} f (x) d x = \int_ {\mathrm{G} / \Gamma} \left(\sum_ {z \in \mathbf {Z} ^ {n}, z \neq 0} f (X \cdot z)\right) d \mu (\dot {X})\tag{1}
\]

\((\dot{X} = X\Gamma)\) . (Let \(G'\) be the subgroup of \(G\) leaving invariant the basis vector \((1,0,0,\ldots,0)\) , so that \(G/G'\) may be identified with \(\mathbf{R}^n\) . Let \(G'' = G' \cap \Gamma\) . Let \(H\) be the subgroup of matrices \(\begin{pmatrix} 1 & w \\ 0 & Y \end{pmatrix}\) of \(G'\) for which \(Y \in \mathbf{SL}(n-1, \mathbf{R})\) has integer entries. Then \(H/G''\) is compact and, by the induction hypothesis, \(G'/H\) has a finite invariant measure, therefore by Prop. 12, §2, No.~8, \(G'/G''\) has a finite invariant measure. Applying Exer. 6 of §2, and Exer. 20 b) of A, VII, §4, show that, for \(f \in \mathcal{K}(\mathbf{R}^n)\) ,

\[
a \int_ {\mathbf {R} ^ {n}} f (x) d x = \int_ {\mathrm{G} / \Gamma} \left(\sum f (X \cdot m)\right) d \mu (\dot {X}),
\]

where the summation is over the set of vectors \(m = (m_1, \ldots, m_n)\) whose coordinates are setwise coprime integers (A, VII, §1, No.~2, Th. 1). Applying this to the functions \(f(2x), f(3x), \ldots\) and summing, prove (1). Then, replacing \(f\) by the function \(x \mapsto \varepsilon^n f(\varepsilon x)\) with \(f \geqslant 0\) , and letting \(\varepsilon\) tend to 0, show that every compact subset of \(G/\Gamma\) has measure \(\leqslant 1\) .

\begin{enumerate}
\def\labelenumi{\arabic{enumi})}
\setcounter{enumi}{7}
\tightlist
\item
  \begin{enumerate}
  \def\labelenumii{\alph{enumii})}
  \tightlist
  \item
    Show that there exists on \(\mathbf{S}_{n-1}\) a measure \(\omega_{n-1}\) invariant under \(\mathbf{O}(n,\mathbf{R})\) and, up to a constant factor, only one (regard \(\mathbf{S}_{n-1}\) as a homogeneous space for \(\mathbf{O}(n,\mathbf{R})\) ). b) The mapping \((t,z)\mapsto tz\) permits identifying the topological space \(\mathbf{R}^n -\{0\}\) with the topological space \(\mathbf{R}_+^*\times \mathbf{S}_{n-1}\) . Show that the measure induced on \(\mathbf{R}^n -\{0\}\) by the Lebesgue measure on \(\mathbf{R}^n\) may then be identified, up to a constant factor, with \(t^{n-1}dt\otimes d\omega_{n-1}(z)\) . (Make use of the fact that the Lebesgue measure on \(\mathbf{R}^n\) is invariant under \(\mathbf{O}(n,\mathbf{R})\) and the fact that the Lebesgue measure of the closed ball with center 0 and radius \(r\) is proportional to \(r^n\) .)
  \end{enumerate}
\end{enumerate}

\begin{enumerate}
\def\labelenumi{\alph{enumi})}
\setcounter{enumi}{2}
\item
  Let \(L_{h}\) be the intersection of \(S_{n-1}\) and the hyperplane \(x_{n}=h\) . The non-empty \(L_{h}\) are the intransitivity classes of \(\mathbf{O}(n-1,\mathbf{R})\) in \(S_{n-1}\) . Show that \(\omega_{n-1}\) may be put in the form \(\int\lambda_{h}d\nu(h)\) , where \(\lambda_{h}\) is a measure invariant under \(\mathbf{O}(n-1,\mathbf{R})\) and carried by \(L_{h}\) , and \(\nu\) is a measure on \([-1,1]\) . Identifying \(L_{h}\) with \(S_{n-2}\) by \(z\mapsto he_{n}+\sqrt{1-h^{2}}z\) , one can assume that \(\lambda_{h}=\omega_{n-2}\) . Let \(K_{\theta}\) be the subset of \(S_{n-1}\) defined by \(\sin\theta\leqslant x_{n}\leqslant1\) . By calculating the Lebesgue measure of the set of tz ( \(0\leqslant t\leqslant1\) , \(z\in K_{\theta}\) ), show that \(\omega_{n-1}(K_{\theta})\) is proportional to \(\int_{\theta}^{\pi/2}\cos^{n-2}\varphi d\varphi\) . From this, deduce that if one sets \(h=\sin\theta(-\pi/2\leqslant\theta\leqslant\pi/2)\) , then the measure \(\nu\) may be identified with \(\cos^{n-2}\theta d\theta\) .
\item
  Conclude, by induction on \(n\) , that
\end{enumerate}

\[
d \omega_ {n - 1} = \cos^ {n - 2} \theta_ {1} \cos^ {n - 3} \theta_ {2} \dots \cos \theta_ {n - 2} d \theta_ {1} d \theta_ {2} \dots d \theta_ {n - 1},
\]

where

\[
\begin{array}{r l} & x _ {1} = \sin \theta_ {1} \\ & x _ {2} = \cos \theta_ {1} \sin \theta_ {2} \\ & \dots \\ & x _ {n - 1} = \cos \theta_ {1} \cos \theta_ {2} \dots \cos \theta_ {n - 2} \sin \theta_ {n - 1} \\ & x _ {n} = \cos \theta_ {1} \cos \theta_ {2} \dots \cos \theta_ {n - 2} \cos \theta_ {n - 1} \end{array}
\]

with \(-\pi / 2 \leqslant \theta_{i} \leqslant \pi / 2\) for \(1 \leqslant i \leqslant n - 2, 0 \leqslant \theta_{n-1} < 2\pi\) .

\begin{enumerate}
\def\labelenumi{\arabic{enumi})}
\setcounter{enumi}{8}
\tightlist
\item
  Let \((e_1, e_2, e_3)\) be the canonical basis of \(\mathbf{R}^3\) . Every matrix \(\sigma \in \mathbf{SO}(3, \mathbf{R})\) such that \(\sigma(e_3) \neq e_3\) and \(\sigma(e_3) \neq -e_3\) may be written in a unique way as
\end{enumerate}

\[
\begin{array}{l} \sigma (\varphi , \psi , \theta) = \\ \left( \begin{array}{c c c} \cos \varphi   \cos \psi - \sin \varphi   \sin \psi   \cos \theta & - \cos \varphi   \sin \psi - \sin \varphi   \cos \psi   \cos \theta & \sin \varphi   \sin \theta \\ \sin \varphi   \cos \psi + \cos \varphi   \sin \psi   \cos \theta & - \sin \varphi   \sin \psi + \cos \varphi   \cos \psi   \cos \theta & - \cos \varphi   \sin \theta \\ \sin \psi   \sin \theta & \cos \psi   \sin \theta & \cos \theta \end{array} \right), \end{array}
\]

where \(0 < \theta < \pi\) , \(0 \leqslant \varphi < 2\pi\) , \(0 \leqslant \psi < 2\pi\) (Euler's angles). Show that \(\sin \theta d\theta d\varphi d\psi\) is a Haar measure on \(\mathbf{SO}(3, \mathbf{R})\) . (Identify the homogeneous space \(\mathbf{SO}(3, \mathbf{R}) / \mathbf{SO}(2, \mathbf{R})\) with \(\mathbf{S}_2\) , use Exer. 8, and Th. 3 of §2.)

\begin{enumerate}
\def\labelenumi{\arabic{enumi})}
\setcounter{enumi}{9}
\tightlist
\item
  Let D be the group of displacements of \(\mathbf{R}^{3}\) with determinant 1, the semi-direct product of \(\mathbf{SO}(3,\mathbf{R})\) and the group T of translations. We employ the notation \(\sigma(\varphi,\psi,\theta)\) of Exer. 9, and denote by \(t(\xi,\eta,\zeta)\) the translation with vector \((\xi,\eta,\zeta)\) .
\end{enumerate}

\begin{enumerate}
\def\labelenumi{\alph{enumi})}
\item
  Let H be the closed subgroup of D that leaves \(\mathbf{R}e_{3}\) stable and preserves orientation on \(\mathbf{R}e_{3}\) . Show that H is the product of the group of translations \(t(0,0,\zeta)\) and the group of rotations \(\sigma(\omega,0,0)\) . From this, deduce that the homogeneous space E = D/H, which may be identified with the space of oriented affine lines of \(\mathbf{R}^{3}\) , possesses a measure invariant under D, and, up to a constant factor, only one. (Make use of Exer. 5 a).
\item
  Let \(\mathbf{E}_1\) be the open subspace of \(\mathbf{E}\) formed by the oriented lines not orthogonal to \(e_3\) ; such an oriented line can be determined by the coordinates \(\xi', \eta'\) of its point of intersection with \(\mathbf{R}e_{1} + \mathbf{R}e_{2}\) , and the coordinates \((\sin \varphi \sin \theta, -\cos \varphi \sin \theta, \cos \theta)\) of a direction vector. Show that a measure on \(E_{1}\) invariant under D is given by
\end{enumerate}

\[
\sin \theta | \cos \theta | d \xi^ {\prime} d \eta^ {\prime} d \varphi d \theta .
\]

\begin{enumerate}
\def\labelenumi{\arabic{enumi})}
\setcounter{enumi}{10}
\tightlist
\item
  Let \(G\) be the compact group \(\mathbf{O}(n, \mathbf{R})\) , \(\lambda\) the normalized Haar measure on \(G\) . For \(0 < k < n\) , let \(H\) be the closed subgroup of \(G\) leaving invariant (globally) the subspace \(\mathbf{R}^k\) of \(\mathbf{R}^n\) ; the homogeneous space \(G / H\) , equipped with its quotient topology, may be identified canonically with the Grassmannian \(E = G_{n-1, k-1}(\mathbf{R})\) (GT, VI, §3, No.~6) of the \(k\) -dimensional linear subspaces of \(\mathbf{R}^n\) . There is a measure \(\mu\) on \(E\) invariant under \(G\) , determined up to a constant. For every subspace \(P \in E\) , let \(\sigma_P\) be the image of the Lebesgue measure of \(\mathbf{R}^k\) under an \(s \in G\) such that \(s \cdot \mathbf{R}^k = P\) (independent of the element \(s \in G\) satisfying this relation). Show that \(\mu\) can be chosen in such a way that the following property is satisfied: for every continuous function \(f\) on \(\mathbf{R}^n\) with compact support, and for every \(P \in E\) , one sets \(F(P) = \int_{P} f(x) d\sigma(x)\) ; then
\end{enumerate}

\[
\int_ {\mathrm{E}} \mathrm{F} (\mathrm{P}) d \mu (\mathrm{P}) = \int_ {\mathbf {R} ^ {n}} \| x \| ^ {k - n} f (x) d x.
\]

(If \(\mathbf{P}_0 = \mathbf{R}^k\) , note that one can write \(\int_{\mathbf{E}} \mathrm{F}(\mathrm{P}) d\mu(\mathrm{P}) = c \int_{\mathbf{G}} \mathrm{F}(s \cdot \mathrm{P}_0) d\lambda(s)\) for a suitable constant \(c\) , and on the other hand,

\[
\int_ {\mathrm{G}} f (s \cdot x) d \lambda (s) = c ^ {\prime} \int_ {\mathbf {S} _ {n - 1}} f (\| x \| z) d \omega_ {n - 1} (z)
\]

with the notations of Exercise 8; finally, make use of Exer. 8 b).)

¶ 12) Let K be a commutative field, E a vector space over K, F a linear subspace of E, p the canonical homomorphism of E onto E/F, and A the set of homomorphisms \(f: E/F \to E\) such that \(p \circ f\) is the identity homomorphism of E/F.

\begin{enumerate}
\def\labelenumi{\alph{enumi})}
\item
  If \(f \in A\) and \(h \in \operatorname{Hom}(\mathrm{E} / \mathrm{F}, \mathrm{F})\) , then \(f + h \in A\) . Show that if \(f, f' \in A\) , there exists one and only one \(h \in \operatorname{Hom}(\mathrm{E} / \mathrm{F}, \mathrm{F})\) such that \(f + h = f'\) . One can therefore regard \(A\) as an affine space of which \(\operatorname{Hom}(\mathrm{E} / \mathrm{F}, \mathrm{F})\) is the space of translations.
\item
  Let B be the set of linear subspaces of E that are supplementary to F in E. Show that \(f \mapsto f(\mathrm{E}/\mathrm{F})\) is a bijection of A onto B. One can therefore regard B as an affine space of which \(\operatorname{Hom}(\mathrm{E}/\mathrm{F}, \mathrm{F})\) is the space of translations.
\item
  Show that if \(u\) is an automorphism of \(E\) such that \(u(F) = F\) , then the mapping \(f \mapsto u \circ f\) is an affine bijection of \(A\) onto \(A\) . (Choose an origin in \(A\) , and observe that the mapping \(h \mapsto u \circ h\) is an automorphism of the vector space \(\operatorname{Hom}(E/F, F)\) .) From this, deduce that the mapping \(F' \mapsto u(F')\) is an affine bijection of \(B\) onto \(B\) .
\item
  Assume that K = R and \(\dim E < +\infty\) . Let G be a compact group, and \(\rho\) a continuous linear representation of G in E such that \(\rho(s)(F) = F\) for all \(s \in G\) . Show that there exists an \(F' \in B\) such that \(\rho(s)(F') = F'\) for all \(s \in G\) . (Make use of c) and Lemma 2 of No.~2.) Obtain this result anew using Prop. 1 of No.~1.
\end{enumerate}

\chapter{Convolution and representations}

\section{CONVOLUTION}

\subsection{Definition and examples}

Recall (Ch. V, §6, Nos. 1 and 4; Ch. VI, §2, No.~10) that, if X and Y are locally compact spaces, \(\mu\) a measure on X, and \(\varphi\) a mapping of X into Y, \(\varphi\) is said to be \(\mu\) -proper if: \(a)\) \(\varphi\) is \(\mu\) -measurable; \(b)\) for every compact subset K of Y, \(\overline{\varphi}^1 (\mathbf{K})\) is essentially \(\mu\) -integrable. Then the image measure \(\nu = \varphi (\mu)\) on Y exists and has the following property: for a function \(\mathbf{f}\) on Y, with values in a Banach space or in \(\overline{\mathbf{R}}\) , to be essentially integrable for \(\nu\) , it is necessary and sufficient that \(\mathbf{f} \circ \varphi\) be so for \(\mu\) , in which case,

\[
\int_ {\mathrm{Y}} \mathbf {f} (y) d \nu (y) = \int_ {\mathrm{X}} \mathbf {f} (\varphi (x)) d \mu (x).
\]

DEFINITION 1. --- Let \(X_1, \ldots, X_n\) be locally compact spaces, \(\mu_i\) a measure on \(X_i\) ( \(1 \leqslant i \leqslant n\) ); let \(X\) be the product of the \(X_i\) , \(\mu\) that of the \(\mu_i\) . Let \(\varphi\) be a mapping of \(X\) into a locally compact space \(Y\) . One says that the sequence \((\mu_i)\) is \(\varphi\) -convolvable, or that \(\mu_1, \ldots, \mu_n\) are \(\varphi\) -convolvable, if \(\varphi\) is \(\mu\) -proper; in this case, the image \(\nu = \varphi(\mu)\) of \(\mu\) under \(\varphi\) is called the convolution product of the \(\mu_i\) for \(\varphi\) , and is denoted \(*^{\varphi(\mu_i)}_{1 \leqslant i \leqslant n}\) , or \(\underset{i=1}{\overset{n}{*}} \mu_i\) , or \(\mu_1 * \mu_2 * \cdots * \mu_n\) .

The last two notations are of course used only when there can be no doubt as to \(\varphi\) .

Let f be a function on Y, with values in a Banach space or in \(\overline{\mathbf{R}}\) . In order that f be essentially integrable for \(\mu_{1}*\cdots*\mu_{n}\) , it is necessary and sufficient that the function

\[
\left(x _ {1}, \dots , x _ {n}\right) \mapsto \mathbf {f} \left(\varphi \left(x _ {1}, \dots , x _ {n}\right)\right)
\]

be essentially integrable for \(\mu_{1}\otimes\mu_{2}\otimes\cdots\otimes\mu_{n}\) , in which case

\[
\int \mathbf {f} d (\mu_ {1} * \dots * \mu_ {n}) = \int \mathbf {f} \big (\varphi (x _ {1}, \ldots , x _ {n}) \big) d \mu_ {1} (x _ {1}) \ldots d \mu_ {n} (x _ {n}),\tag{1}
\]

a formula that may be regarded as defining \(\mu_{1}*\cdots *\mu_{n}\) when one takes \(\mathbf{f}\in\mathcal{K}(\mathrm{Y})\) .

The definitions imply at once that the \(\mu_{i}\) are convolvable if and only if the \(|\mu_{i}|\) are. When this is the case,

\[
\left| \varphi \left(\mu_ {1} \otimes \dots \otimes \mu_ {n}\right) \right| \leqslant \varphi \left(\left| \mu_ {1} \otimes \dots \otimes \mu_ {n} \right|\right) = \varphi \left(\left| \mu_ {1} \right| \otimes \dots \otimes \left| \mu_ {n} \right|\right)
\]

(Ch. VI, §2, No.~10), that is,

\[
\left| \begin{array}{c} * \\ i \end{array} \mu_ {i} \right| \leqslant * _ {i} | \mu_ {i} |.\tag{2}
\]

If the \(\mu_{i}\) are convolvable and positive, and if \(\nu_{i}\) is a measure on \(X_{i}\) such that \(0 \leqslant \nu_{i} \leqslant \mu_{i}\) , then the \(\nu_{i}\) are convolvable and

\[
\underset {i} {\ast} \nu_ {i} \leqslant \underset {i} {\ast} \mu_ {i}.
\]

Suppose \(\mu_1, \mu_2, \ldots, \mu_n\) are convolvable, and that \(\mu_1', \mu_2, \ldots, \mu_n\) are convolvable ( \(\mu_1'\) being a measure on \(X_1\) ). By Ch. V, §6, No.~3, Cor. 1 of Prop. 6, \(\mu_1 + \mu_1', \mu_2, \ldots, \mu_n\) are convolvable and

\[
\left(\mu_ {1} + \mu_ {1} ^ {\prime}\right) * \mu_ {2} * \dots * \mu_ {n} = \mu_ {1} * \mu_ {2} * \dots * \mu_ {n} + \mu_ {1} ^ {\prime} * \mu_ {2} * \dots * \mu_ {n}.
\]

Examples. --- 1) For any \(\varphi\) , the measures \(\varepsilon_{x_i}\) , where \(x_i \in X_i\) for \(1 \leqslant i \leqslant n\) , are always convolvable and have convolution product \(\varepsilon_y\) , with \(y = \varphi(x_1, x_2, \ldots, x_n)\) . Consequently, if each of the \(\mu_i\) has finite support, then the \(\mu_i\) are convolvable and \(\mu_1 * \cdots * \mu_n\) has finite support. In particular, let M be a monoid \(^1\) equipped with a locally compact topology; if one takes \(\varphi\) to be the law of composition in M then the measures on M with finite support form, for convolution, an algebra that is none other than the algebra of the monoid M (over R or over C, according as one considers real or complex measures) (A, III, §2, No.~6).

\begin{enumerate}
\def\labelenumi{\arabic{enumi})}
\setcounter{enumi}{1}
\tightlist
\item
  Let M be a monoid equipped with the discrete topology; assume that for each \(m \in M\) , there are only finitely many pairs \((m', m'') \in M \times M\) such that \(m'm'' = m\) ; this amounts to saying that the law of composition in M is a proper mapping of \(M \times M\) into M; the measures on M then form an algebra for convolution, an algebra that is none other than the total algebra of the monoid M (A, III, §2, No.~10); we note the following two special cases:
\end{enumerate}

\begin{enumerate}
\def\labelenumi{\alph{enumi})}
\tightlist
\item
  M = N, the law of composition being addition. To every measure \(\mu\) on N, let us associate the formal series
\end{enumerate}

\[
\mathrm{S} (\mu) = \sum_ {n = 0} ^ {\infty} \mu (\{n \}) t ^ {n}
\]

in an indeterminate t. Then \(S(\mu * \mu') = S(\mu)S(\mu')\) . An analogous remark holds for formal series in any number of indeterminates.

\(^{*}b)\) M = N \(^{*}\) , the law of composition being multiplication. To every measure \(\mu\) on N \(^{*}\) , let us associate the formal Dirichlet series

\[
\mathrm{D} (\mu) = \sum_ {n = 1} ^ {\infty} \mu (\{n \}) n ^ {- s}.
\]

Then \(\mathrm{D}(\mu * \mu') = \mathrm{D}(\mu) \mathrm{D}(\mu')\) .

\begin{enumerate}
\def\labelenumi{\arabic{enumi})}
\setcounter{enumi}{2}
\tightlist
\item
  Let X,Y,Z be locally compact spaces, \(\varphi\) a continuous mapping of \(X \times Y\) into Z. If \(x \in X\) and \(\mu\) is a measure on Y, to say that \(\varepsilon_{x}\) and \(\mu\) are \(\varphi\) -convolvable comes to saying that the mapping \(\varphi(x,\cdot)\) of Y into Z is \(\mu\) -proper. One then has \(\varepsilon_{x} * \mu = \varphi(x,\cdot)(\mu)\) .
\end{enumerate}

\subsection{Associativity}

The following lemma completes Prop. 11 of Ch. V, §8, No.~5: \(^{2}\)

Lemma 1. --- For \(1 \leqslant i \leqslant n\) , let \(\mathrm{X}_i, \mathrm{Y}_i\) be two locally compact spaces, \(\mu_i\) a measure on \(\mathrm{X}_i\) , and \(\varphi_i\) a continuous mapping of \(\mathrm{X}_i\) into \(\mathrm{Y}_i\) . Let \(\mathrm{X} = \prod_i \mathrm{X}_i\) , \(\mathrm{Y} = \prod_i \mathrm{Y}_i\) , \(\mu = \bigotimes_i \mu_i\) , and \(\varphi\) the mapping of \(\mathrm{X}\) into \(\mathrm{Y}\) that is the product of the \(\varphi_i\) . If \(\varphi\) is \(\mu\) -proper and \(\mu_i \neq 0\) for each \(i\) , then the \(\varphi_i\) are \(\mu_i\) -proper and \(\varphi(\mu) = \bigotimes_i \varphi_i(\mu_i)\) .

We can suppose that the \(\mu_{i}\) are positive and \(n = 2\) . Let \(f_{1} \in \mathcal{K}_{+}(Y_{1})\) . Since \(\mu_{2} \neq 0\) , there exists an \(f_{2} \in \mathcal{K}_{+}(Y_{2})\) such that \(f_{2} \circ \varphi_{2}\) is not \(\mu_{2}\) -negligible. The function \((x_{1}, x_{2}) \mapsto f_{1}(\varphi_{1}(x_{1})) f_{2}(\varphi_{2}(x_{2}))\) is essentially \(\mu\) -integrable and continuous, hence \(\mu\) -integrable. Therefore there exists an \(x_{2} \in X_{2}\) such that \(f_{2}(\varphi_{2}(x_{2})) \neq 0\) and such that the function \(x_{1} \mapsto f_{1}(\varphi_{1}(x_{1})) f_{2}(\varphi_{2}(x_{2}))\) is \(\mu_{1}\) -integrable. Therefore \(f_{1} \circ \varphi_{1}\) is \(\mu_{1}\) -integrable, which proves that \(\varphi_{1}\) is \(\mu_{1}\) -proper. One argues similarly for \(\varphi_{2}\) . Then \(\varphi(\mu) = \bigotimes_{i} \varphi_{i}(\mu_{i})\) by Prop. 11 of Ch. V, §8, No.~5.

The following lemma completes Prop. 4 of Ch. V, §6, No.~3. \(^{3}\)

Lemma 2. --- Let \(\mathrm{T}, \mathrm{T}', \mathrm{T}''\) be three locally compact spaces, \(\mu\) a measure on \(\mathrm{T}\) , \(\pi\) a \(\mu\) -measurable mapping of \(\mathrm{T}\) into \(\mathrm{T}'\) , \(\pi'\) a continuous mapping of \(\mathrm{T}'\) into \(\mathrm{T}''\) , and \(\pi'' = \pi' \circ \pi\) . If \(\pi''\) is \(\mu\) -proper, then \(\pi\) is \(\mu\) -proper, \(\pi'\) is \(\pi(\mu)\) -proper, and \(\pi''(\mu) = \pi' (\pi(\mu))\) .

Let \(K'\) be a compact subset of \(T'\) . Then \(K'' = \pi'(K')\) is compact, therefore \(\pi''(K'')\) is essentially \(\mu\) -integrable, therefore \(\overline{\pi}(K') \subset \overline{\pi''}(K'')\) is essentially \(\mu\) -integrable, thus \(\pi\) is \(\mu\) -proper. Then \(\pi'\) is \(\pi(\mu)\) -proper and \(\pi''(\mu) = \pi'(\pi(\mu))\) by Ch. V, §6, No.~3, Prop. 4.

PROPOSITION 1. --- Let \(\mathrm{X}_{ij}\) ( \(1 \leqslant i \leqslant m, 1 \leqslant j \leqslant n_i\) ), \(\mathrm{Y}_i\) ( \(1 \leqslant i \leqslant m\) ), and \(\mathbf{Z}\) be locally compact spaces; for each \(i\) , let \(\varphi_i\) be a mapping of \(\mathrm{X}_i = \prod_j \mathrm{X}_{ij}\) into \(\mathrm{Y}_i\) ; let \(\varphi\) be the product of the \(\varphi_i\) , mapping \(\mathrm{X} = \prod_i \mathrm{X}_i\) into \(\mathrm{Y} = \prod_j \mathrm{Y}_i\) ; let \(\psi\) be a mapping of \(\mathrm{Y}\) into \(\mathbf{Z}\) .

\begin{enumerate}
\def\labelenumi{(\roman{enumi})}
\tightlist
\item
  Let \(\mu_{ij}\) be measures given, respectively, on the \(\mathrm{X}_{ij}\) , such that, for each \(i\) , the \(\mu_{ij}\) ( \(1 \leqslant j \leqslant n_i\) ) are \(\varphi_i\) -convolvable, and such that the measures \(*|_{\mu_{ij}}|\) are \(\psi\) -convolvable; then the \(\mu_{ij}\) , for \(1 \leqslant i \leqslant m\) , \(1 \leqslant j \leqslant n_i\) , are \((\psi \circ \varphi)\) -convolvable and
\end{enumerate}

\[
\underset {i, j} {\ast} \mu_ {i j} = \underset {i} {\ast} \left(\underset {j} {\ast} \mu_ {i j}\right).\tag{3}
\]

\begin{enumerate}
\def\labelenumi{(\roman{enumi})}
\setcounter{enumi}{1}
\tightlist
\item
  Assume \(\psi\) and the \(\varphi_{i}\) continuous, and let \(\mu_{ij}\) be measures \(\neq 0\) given, respectively, on the \(X_{ij}\) and \((\psi \circ \varphi)\) -convolvable; then, for each \(i\) , the \(\mu_{ij}\) ( \(1 \leqslant j \leqslant n_{i}\) ) are \(\varphi_{i}\) -convolvable, the measures \(*_{j} |\mu_{ij}|\) are \(\psi\) -convolvable, and the formula (3) holds.
\end{enumerate}

It suffices to consider the case that all of the measures in question are \(\geqslant 0\) .

Let us place ourselves under the hypotheses of (i). The mapping \(\varphi\) is proper for \(\bigotimes_{i,j} \mu_{ij}\) , and

\[
\varphi \left(\bigotimes_ {i, j} \mu_ {i j}\right) = \bigotimes_ {i} \varphi_ {i} \left(\bigotimes_ {j} \mu_ {i j}\right) = \bigotimes_ {i} \left(* _ {j} \mu_ {i j}\right)
\]

(Ch. V, §8, No.~5, Prop. 11). The mapping \(\psi \circ \varphi\) is proper for \(\bigotimes_{i,j} \mu_{ij}\) , and

\[
(\psi \circ \varphi) \left(\bigotimes_ {i, j} \mu_ {i j}\right) = \psi \left(\bigotimes_ {i} \left(* _ {j} \mu_ {i j}\right)\right) = * _ {i} \left(* _ {j} \mu_ {i j}\right)
\]

(Ch. V, §6, Prop. 4). Therefore the \(\mu_{ij}\) ( \(1 \leqslant i \leqslant m\) , \(1 \leqslant j \leqslant n_i\) ) are \((\psi \circ \varphi)\) -convolvable and formula (3) holds.

Let us place ourselves under the hypotheses of (ii). First of all, Lemma 2 proves that \(\varphi\) is proper for \(\bigotimes_{i,j} \mu_{ij}\) . Lemma 1 then proves that for every \(i\) , \(\varphi_i\) is proper for \(\bigotimes_{i} \mu_{ij}\) , and that

\[
\varphi \Big (\bigotimes_ {i, j} \mu_ {i j} \Big) = \bigotimes_ {i} \Big (* _ {j} \mu_ {i j} \Big).
\]

By Lemma 2, \(\psi\) is proper for \(\bigotimes_{i}\left(*_{j}\mu_{ij}\right)\) . Whence the proposition.

COROLLARY. --- Let \(X_{i}, X_{i}^{\prime}\) ( \(1 \leqslant i \leqslant n\) ), Y, \(Y^{\prime}\) be locally compact spaces; let \(\varphi, \varphi^{\prime}\) be continuous mappings of \(X = \prod_{i} X_{i}\) into Y and of \(X^{\prime} = \prod_{i} X_{i}^{\prime}\) into \(Y^{\prime}\) , respectively; let \(f_{i}\) be continuous mappings of \(X_{i}\) into \(X_{i}^{\prime}\) ( \(1 \leqslant i \leqslant n\) ) and g a continuous mapping of Y into \(Y^{\prime}\) , such that \(\varphi^{\prime} \circ f = g \circ \varphi\) , f being the mapping of X into \(X^{\prime}\) that is the product of the \(f_{i}\) . Let \(\mu_{i}\) be measures given respectively on the \(X_{i}\) , all \(\neq 0\) . Then the following two assertions are equivalent:

\begin{enumerate}
\def\labelenumi{(\roman{enumi})}
\item
  \(f_{i}\) is \(\mu_{i}\) -proper for all \(i\) , and the measures \(f_{i}(|\mu_{i}|)\) are \(\varphi'\) -convolvable;
\item
  the \(\mu_{i}\) are \(\varphi\) -convolvable, and \(g\) is proper for \(*_{\varphi}(|\mu_i|)\) .
\end{enumerate}

Moreover, when these assertions are verified,

\[
\left. \ast_ {\varphi^ {\prime}} \left(f _ {i} \left(\mu_ {i}\right)\right) = g \left(\ast_ {\varphi} \left(\mu_ {i}\right)\right) = \ast_ {g \circ \varphi} \left(\mu_ {i}\right). \right.\tag{4}
\]

For, let \(h = \varphi' \circ f = g \circ \varphi\) . By Prop. 1, the conditions (i) and (ii) are each equivalent to the following condition:

\begin{enumerate}
\def\labelenumi{(\roman{enumi})}
\setcounter{enumi}{2}
\tightlist
\item
  the \(\mu_{i}\) are \(h\) -convolvable.
\end{enumerate}

When this is so,

\[
\left. \ast_ {\varphi^ {\prime}} \left(f _ {i} \left(\mu_ {i}\right)\right) = \ast_ {h} \left(\mu_ {i}\right) = g \left(\ast_ {\varphi} \left(\mu_ {i}\right)\right). \right.
\]

\subsection{The case of bounded measures}

PROPOSITION 2. --- Let \(X_{1},\ldots,X_{n},Y\) be locally compact spaces, \(\mu_{i}\) a bounded measure on \(X_{i}\) ( \(1\leqslant i\leqslant n\) ), \(\mu\) the product of the \(\mu_{i}\) , \(\varphi\) a \(\mu\) -measurable mapping of \(\prod_{i}X_{i}\) into Y. Then the \(\mu_{i}\) are \(\varphi\) -convolvable and

\[
\left\| \begin{array}{c} n \\ * \\ i = 1 \end{array} \mu_ {i} \right\| \leqslant \prod_ {i = 1} ^ {n} \| \mu_ {i} \|.
\]

If the \(\mu_{i}\) are moreover positive, then \(\left\| \begin{array}{c}n\\ *\\ i = 1 \end{array} \mu_i\right\| = \prod_{i = 1}^{n}\| \mu_{i}\|\)

For, \(\mu_{i}^{\prime}=|\mu_{i}|\) is bounded and \(\|\mu_{i}^{\prime}\|=\|\mu_{i}\|\) (Ch. III, §1, No.~8, Cor. 1 of Prop. 10). One has \(|\mu_{1}\otimes\cdots\otimes\mu_{n}|=\mu_{1}^{\prime}\otimes\cdots\otimes\mu_{n}^{\prime}\) (Ch. III, §4, Nos. 2, 4), therefore \(\mu_{1}\otimes\cdots\otimes\mu_{n}\) is bounded and

\[
\left\| \mu_ {1} \otimes \dots \otimes \mu_ {n} \right\| = \left\| \mu_ {1} \right\| \dots \left\| \mu_ {n} \right\|
\]

(ibid., Prop. 4). Therefore \(\varphi\) is \(\mu\) -proper (Ch. V, §6, No.~1, Remark 1), that is, the \(\mu_i\) are \(\varphi\) -convolvable. One has \(\left\| \begin{array}{c} n \\ * \\ 1 \end{array} \mu_i'\right\| = \| \mu_1' \otimes \cdots \otimes \mu_n'\|\) (Ch. V, §6, No.~2, Th. 1), consequently \(\left\| \begin{array}{c} n \\ * \\ i=1 \end{array} \mu_i'\right\| = \| \mu_1'\| \cdots \| \mu_n'\|\) . Finally, \(| *_{i} \mu_i| \leqslant *_{i} \mu_i'\) (No.~1, formula (2)), therefore

\[
\left\| \ast_ {i} \mu_ {i} \right\| \leqslant \left\| \ast_ {i} \mu_ {i} ^ {\prime} \right\| = \prod_ {i = 1} ^ {n} \| \mu_ {i} \|.
\]

PROPOSITION 3. --- Let \(X_1, \ldots, X_n, Y\) be locally compact spaces, \(\varphi\) a continuous mapping of \(\prod_{i=1}^{n} X_i\) into \(Y\) . Then the mapping

\[
(\mu , \dots , \mu_ {n}) \mapsto * _ {\varphi} (\mu_ {i})
\]

of \(\prod_{i=1}^{n} \mathcal{M}^1(\mathrm{X}_i)\) into \(\mathcal{M}^1(\mathrm{Y})\) is a continuous multilinear mapping.

This follows from Prop. 2 and what has been said in No.~1.

\subsection{Properties concerning supports}

PROPOSITION 4. --- Let \(X_{1},\ldots,X_{n},Y\) be locally compact spaces, \(\mu_{i}\) a measure on \(X_{i}\) ( \(1\leqslant i\leqslant n\) ), \(S_{i}\) its support, and \(\varphi\) a continuous mapping of \(\prod_{i} X_{i}\) into Y such that the restriction of \(\varphi\) to \(\prod_{i} S_{i}\) is proper. Then the \(\mu_{i}\) are \(\varphi\) -convolvable.

For, let K be a compact subset of Y. The support of \(\mu = \mu_{1} \otimes \cdots \otimes \mu_{n}\) is \(S = \prod_{i} S_{i}\) (Ch. III, §4, No.~2, Prop. 2). Therefore \(\overline{\varphi}(K) \cap \left( \prod_{i} X_{i} - S \right)\) is \(\mu\) -negligible. On the other hand, \(\overline{\varphi}(K) \cap S\) is compact. Therefore \(\overline{\varphi}(K)\) is \(\mu\) -integrable.

PROPOSITION 5. --- Let \(X_1, \ldots, X_n, Y\) be locally compact spaces, \(\mu_i\) a measure on \(X_i\) ( \(1 \leqslant i \leqslant n\) ), \(\mu\) the product of the \(\mu_i\) , \(\varphi\) a \(\mu\) -proper mapping of \(\prod_i X_i\) into \(Y\) , and \(S_i\) the support of \(\mu_i\) .

\begin{enumerate}
\def\labelenumi{\alph{enumi})}
\item
  The support of \(*\mu_i\) is contained in the closure of \(\varphi\left(\prod_{i}S_i\right)\) .
\item
  If \(\varphi\) is continuous and the \(\mu_{i}\) are positive, then the support of \(*_{i} \mu_{i}\) is the closure of \(\varphi\left(\prod_{i} S_{i}\right)\) .
\end{enumerate}

Let \(S = \prod_{i} S_{i}\) be the support of \(\mu\) . The support of \(*_{i} \mu_{i}\) is contained in \(\overline{\varphi(S)}\) by Ch. V, §6, No.~2, Cor. 3 of Prop. 2. If \(\varphi\) is continuous and the \(\mu_{i}\) are positive, then the support of \(*_{i} \mu_{i}\) is \(\overline{\varphi(S)}\) (loc. cit., Cor. 4 of Prop. 2).

COROLLARY. --- If \(\varphi\) is continuous and the \(\mu_{i}\) have compact support, then the \(\mu_{i}\) are convolvable and \(*\mu_{i}\) has compact support.

\subsection{Vectorial expression of the convolution product}

PROPOSITION 6. --- Let X,Y,Z be locally compact spaces, \(\varphi\) a continuous mapping of \(X \times Y\) into Z, and \(\lambda, \mu\) measures on X,Y. For \(\lambda\) and \(\mu\) to be \(\varphi\) -convolvable, it is necessary and sufficient that the mapping \((x,y) \mapsto \varepsilon_{\varphi(x,y)} = \varepsilon_x * \varepsilon_y\) of \(X \times Y\) into \(\mathcal{M}(Z)\) be scalarly \((\lambda \otimes \mu)\) -integrable for the topology \(\sigma(\mathcal{M}(Z), \mathcal{K}(Z))\) , in which case

\[
\lambda * \mu = \int_ {\mathrm{X} \times \mathrm{Y}} \left(\varepsilon_ {x} * \varepsilon_ {y}\right) d \lambda (x) d \mu (y).
\]

To say that \(\lambda\) and \(\mu\) are \(\varphi\) -convolvable signifies that, for every \(f \in \mathcal{H}(\mathbf{Z})\) , \(f \circ \varphi\) is \((\lambda \otimes \mu)\) -integrable, that is, for every \(f \in \mathcal{H}(\mathbf{Z})\) the function \((x,y) \mapsto \langle f, \varepsilon_{\varphi(x,y)} \rangle\) is \((\lambda \otimes \mu)\) -integrable, that is, again, that the mapping \((x,y) \mapsto \varepsilon_{\varphi(x,y)}\) of \(\mathrm{X} \times \mathrm{Y}\) into \(\mathcal{M}(\mathbf{Z})\) is scalarly \((\lambda \otimes \mu)\) -integrable for \(\sigma (\mathcal{M}(\mathbb{Z}),\mathcal{K}(\mathbb{Z}))\) . If this is the case, then

\[
\langle \lambda * \mu , f \rangle = \int f (\varphi (x, y)) d \lambda (x) d \mu (y) = \int_ {\mathrm{X} \times \mathrm{Y}} \left\langle \varepsilon_ {\varphi (x, y)}, f \right\rangle d \lambda (x) d \mu (y),
\]

whence \(\lambda * \mu = \int_{\mathbf{X} \times \mathbf{Y}} \varepsilon_{\varphi(x, y)} d\lambda(x) d\mu(y)\) .

PROPOSITION 7. --- Let X,Y,Z be locally compact spaces, \(\varphi\) a continuous mapping of \(X \times Y\) into Z, and \(\lambda, \mu\) measures on X,Y. Assume that for every \(x \in X\) , \(\varepsilon_x\) and \(\mu\) are \(\varphi\) -convolvable. For \(\lambda\) and \(\mu\) to be \(\varphi\) -convolvable, it is necessary and sufficient that the mapping \(x \mapsto \varepsilon_x * |\mu|\) of X into \(\mathcal{M}(Z)\) be scalarly \(\lambda\) -integrable for the topology \(\sigma(\mathcal{M}(Z), \mathcal{K}(Z))\) , in which case \(\lambda * \mu = \int_{X} (\varepsilon_x * \mu) d\lambda(x)\) .

Suppose that \(\lambda\) and \(\mu\) are \(\varphi\) -convolvable. For every \(f \in \mathcal{K}(\mathbb{Z})\) , \(f \circ \varphi\) is \((|\lambda| \otimes |\mu|)\) -integrable, therefore the function \(x \mapsto \int_{\mathbb{Y}} f(\varphi(x, y)) d|\mu|(y) = \langle f, \varepsilon_x * |\mu| \rangle\) (which by hypothesis is defined for all \(x \in \mathbb{X}\) ) is \(\lambda\) -integrable; thus \(x \mapsto \varepsilon_x * |\mu|\) is scalarly \(\lambda\) -integrable for \(\sigma(\mathcal{M}(\mathbb{Z}), \mathcal{K}(\mathbb{Z}))\) , and

\[
\langle f, \lambda * \mu \rangle = \int_ {\mathrm{X}} d \lambda (x) \int_ {\mathrm{Y}} f (\varphi (x, y)) d \mu (y) = \int_ {\mathrm{X}} \left\langle f, \varepsilon_ {x} * \mu \right\rangle d \lambda (x),
\]

whence \(\lambda * \mu = \int_{\mathbf{X}} (\varepsilon_x * \mu) d\lambda(x)\) . Conversely, suppose that the mapping \(x \mapsto \varepsilon_x * |\mu|\) of X into \(\mathcal{M}(\mathbf{Z})\) is scalarly \(\lambda\) -integrable for \(\sigma(\mathcal{M}(\mathbf{Z}), \mathcal{K}(\mathbf{Z}))\) . Let \(f \in \mathcal{H}_+(\mathbf{Z})\) . Then the function \((x,y) \mapsto f(\varphi(x,y))\) is continuous and (Ch. V, §8, No.~3, Prop. 5)

\[
\begin{array}{r l} \iint^ {*} f (\varphi (x, y)) d | \lambda | (x) d | \mu | (y) & = \int^ {*} d | \lambda | (x) \int^ {*} f (\varphi (x, y)) d | \mu | (y) \\ & = \int^ {*} \langle f, \varepsilon_ {x} * | \mu | \rangle d | \lambda | (x) <   + \infty . \end{array}
\]

Therefore \(f \circ \varphi\) is \((\lambda \otimes \mu)\) -integrable, so that \(\lambda\) and \(\mu\) are \(\varphi\) -convolvable.

\section{LINEAR REPRESENTATIONS OF GROUPS}

\subsection{Continuous linear representations}

Let \(\mathbf{G}\) be a topological group, \(\mathbf{E}\) a locally convex space, \(U\) a linear representation of \(\mathbf{G}\) in \(\mathbf{E}\) .

DEFINITION 1. --- (i) U is said to be separately continuous if, for every \(s \in G\) , \(U(s)\) is a continuous endomorphism of E, and if, for every \(x \in E\) , the mapping \(s \mapsto U(s)x\) of G into E is continuous.

\begin{enumerate}
\def\labelenumi{(\roman{enumi})}
\setcounter{enumi}{1}
\item
  \(U\) is said to be continuous if \((s,x)\mapsto U(s)x\) is a continuous mapping of \(\mathbf{G}\times \mathbf{E}\) into \(\mathbf{E}\) .
\item
  \(U\) is said to be equicontinuous if it is continuous and if the set of endomorphisms \(U(s)\) , where \(s\) runs over \(G\) , is equicontinuous.
\end{enumerate}

Remarks. --- 1) To say that U is separately continuous means that \(s \mapsto U(s)\) is a continuous mapping of G into the space \(\mathcal{L}(\mathrm{E};\mathrm{E})\) of continuous endomorphisms of E, equipped with the topology of pointwise convergence.

\begin{enumerate}
\def\labelenumi{\arabic{enumi})}
\setcounter{enumi}{1}
\tightlist
\item
  To say that \(U\) is continuous is equivalent to the following set of three conditions:
\end{enumerate}

\begin{enumerate}
\def\labelenumi{\alph{enumi})}
\tightlist
\item
  for every \(s \in G\) , \(U(s)\) is continuous; b) there exists a neighborhood V of e such that \(U(V)\) is equicontinuous; c) there exists a total set D in E such that, for every \(x \in D\) , the mapping \(s \mapsto U(s)x\) is continuous.
\end{enumerate}

These conditions are obviously necessary. Conversely, suppose that the conditions \(a\) , \(b\) , \(c\) are satisfied. On \(U(\mathrm{V})\) , the topology of pointwise convergence is identical to the topology of pointwise convergence in D (TVS, III, §3, No.~4, Prop. 5). Therefore the mapping \((s,x)\mapsto U(s)x\) of \(\mathbf{V}\times \mathbf{E}\) into E is continuous (GT, X, §2, No.~1, Cor. 3 of Prop. 1). Since \(U(s_0s)x = U(s_0)(U(s)x)\) for all \(s_0\in \mathbf{G}\) , \(s\in \mathbf{G}\) , \(x\in \mathbf{E}\) , one sees that \(U\) is continuous.

When \(\mathbf{G}\) is locally compact, the conditions \(a)\) and \(b)\) are equivalent to the condition:

\(a^{\prime})\) for every compact subset \(\mathbf{K}\) of \(\mathbf{G}\) , \(U(\mathbf{K})\) is equicontinuous.

\begin{enumerate}
\def\labelenumi{\arabic{enumi})}
\setcounter{enumi}{2}
\item
  Suppose that \(U\) is a continuous linear representation of \(G\) in \(E\) . For every \(s \in G\) , let \(\widehat{U}(s)\) be the continuous extension of \(U(s)\) to the completion \(\widehat{E}\) of \(E\) . Then \(\widehat{U}\) is a linear representation of \(G\) in \(\widehat{E}\) , satisfying conditions \(a)\) and \(c)\) of Remark 2, and also condition \(b)\) by GT, X, §2, No.~2, Prop. 4. Therefore \(\widehat{U}\) is a continuous linear representation of \(G\) in \(\widehat{E}\) .
\item
  When \(\mathbf{E}\) is a normed space, \(U\) is said to be isometric if \(\| U(s)\| = 1\) for every \(s\in \mathbf{G}\) . For this, it suffices that \(\| U(s)\| \leqslant 1\) for all \(s\in \mathbf{G}\) , because one then has
\end{enumerate}

\[
1 = \| 1 \| \leqslant \| U (s) \| \cdot \| U (s ^ {- 1}) \|,
\]

whence \(\| U(s)\| = \| U(s^{-1})\| = 1\) for all \(s\in G\) .

PROPOSITION 1. --- If G is a locally compact group and E is barreled, then every separately continuous linear representation U of G in E is continuous.

For every compact subset K of G, \(U(\mathrm{K})\) is compact for the topology of pointwise convergence (Remark 1), therefore is equicontinuous (TVS, III, §4, No.~2, Th. 1); one then applies Remark 2.

Lemma 1. --- Let G be a locally compact group, \(\rho\) a lower semi-continuous finite numerical function \(\geqslant 0\) on G such that \(\rho(st) \leqslant \rho(s)\rho(t)\) for all \(s, t \in \mathbf{G}\) . Then \(\rho\) is bounded above on every compact subset of G.

There exists a nonempty open subset U of G such that \(\rho\) is bounded above on U (GT, IX, §5, No.~4, Th. 2). Let K be a compact subset of G. Then K is covered by a finite number of sets \(s_{1}U,\ldots,s_{n}U\) . For every \(x\in U\) , one has \(\rho(s_{i}x)\leqslant\rho(s_{i})\rho(x)\) , therefore \(\rho\) is bounded above on the \(s_{i}U\) , hence on K.

Lemma 2. --- Let G be a topological group, U a linear representation of G in a normed space E, and A a dense subset of E. Assume that for every \(s \in G\) , \(U(s)\) is continuous, and that, for every \(x \in A\) , \(s \mapsto U(s)x\) is a continuous mapping of G into E. Then the function \(s \mapsto g(s) = \|U(s)\|\) on G is lower semi-continuous and satisfies \(g(st) \leqslant g(s)g(t)\) .

Let B be the unit ball of E. Then \(g(s) = \sup_{x \in \mathbf{B} \cap \mathbf{A}} \|U(s)x\|\) , and each function \(s \mapsto \|U(s)x\|\) is continuous on G, therefore g is lower semicontinuous. On the other hand,

\[
g (s t) = \| U (s) U (t) \| \leqslant \| U (s) \| \cdot \| U (t) \| = g (s) g (t).
\]

PROPOSITION 2. --- Let G be a locally compact group, U a linear representation of G in a normed space E. Let A be a dense subset of E. Assume that for every \(s \in G\) , \(U(s)\) is continuous and that, for every \(x \in A\) , \(s \mapsto U(s)x\) is a continuous mapping of G into E. Then U is continuous.

For, \(\|U(s)\|\) is bounded on every compact subset of G by Lemmas 1 and 2, and one then applies Remark 2.

\subsection{Contragredient representation}

Let U be a separately continuous linear representation of G in E. Let \(E'\) be the dual of E. The mapping \(s \mapsto {}^{t}U(s)\) is a linear representation in \(E'\) of the group \(G^{0}\) opposite G; we shall say that this representation is the transpose of U. The mapping \(s \mapsto {}^{t}U(s^{-1}) = {}^{t}U(s)^{-1}\) is a linear representation of G in \(E'\) , called the contragredient of U.

Lemma 3. --- Let X be a locally compact space, Y and Z topological spaces, \(\varphi\) a continuous mapping of \(X \times Y\) into \(Z\) , and \(\varphi_x\) the mapping \(y \mapsto \varphi(x, y)\) of Y into \(Z\) . The spaces \(\mathcal{C}(Y), \mathcal{C}(Z)\) being equipped with the topology of compact convergence, the mapping \((x, f) \mapsto f \circ \varphi_x\) of \(X \times \mathcal{C}(Z)\) into \(\mathcal{C}(Y)\) is continuous.

It clearly suffices to consider the case that \(\mathbf{X}\) is compact. Let \((x_0, f_0) \in \mathbf{X} \times \mathcal{C}(\mathbf{Z})\) , \(\mathbf{K}\) a compact subset of \(\mathbf{Y}\) , and \(\varepsilon > 0\) . Let \(\mathbf{K}' = \varphi(\mathbf{X} \times \mathbf{K})\) . Since \(f_0 \circ \varphi\) is uniformly continuous in \(X \times K\) , there exists a neighborhood \(W\) of \(x_0\) such that \(|f_0(\varphi(x, y)) - f_0(\varphi(x_0, y))| \leqslant \varepsilon\) for \(x \in W\) and \(y \in K\) . On the other hand, if one takes \(f \in \mathcal{C}(Z)\) to be such that \(|f(z) - f_0(z)| \leqslant \varepsilon\) for all \(z \in K'\) , one will have \(|f(\varphi(x, y)) - f_0(\varphi(x, y))| \leqslant \varepsilon\) for \(x \in X\) , \(y \in K\) , and therefore \(|f(\varphi(x, y)) - f_0(\varphi(x_0, y))| \leqslant 2\varepsilon\) for \(x \in W\) , \(y \in K\) . Whence the lemma.

Let us now return to the earlier notations.

PROPOSITION 3. --- (i) If \(U\) is separately continuous, then \({}^t U\) is separately continuous when \(\mathbf{E}'\) is equipped with the weak topology \(\sigma(\mathbf{E}',\mathbf{E})\) .

\begin{enumerate}
\def\labelenumi{(\roman{enumi})}
\setcounter{enumi}{1}
\tightlist
\item
  If \(\mathbf{G}\) is locally compact and \(U\) is continuous, then \({}^t U\) is continuous when \(\mathbf{E}'\) is equipped with the topology of compact convergence.
\end{enumerate}

The assertion (i) is immediate. The assertion (ii) follows from Lemma 3 where one has taken \(\mathbf{X} = \mathbf{G}\) , \(\mathbf{Y} = \mathbf{Z} = \mathbf{E}\) , \(\varphi(s, x) = U(s)x\) .

\subsection{Example: linear representations in spaces of continuous functions}

Let G be a discrete group operating on the left on a set X. A complex function \(\chi\) on \(G \times X\) is called a multiplier if

\begin{enumerate}
\def\labelenumi{(\arabic{enumi})}
\tightlist
\item
\end{enumerate}

\[
\chi (e, x) = 1 \quad \text { for   all } x \in X;\tag{2}
\]

\[
\chi (s t, x) = \chi (s, t x) \chi (t, x) \quad \text { for   all } s, t \text { in } G, x \in X.
\]

It follows that

\[
\chi (t ^ {- 1}, t x) \chi (t, x) = 1 \quad \text { for   all } t \in \mathrm{G}, x \in \mathrm{X},\tag{3}
\]

and in particular \(\chi(t, x) \neq 0\) for all \(t \in \mathbf{G}\) , \(x \in \mathbf{X}\) .

For every complex function \(f\) defined on \(\mathbf{X}\) and every \(s \in \mathbf{G}\) , let \(\pmb{\gamma}_{\chi}(s)f\) be the complex function on \(\mathbf{X}\) defined by

\[
\left(\boldsymbol {\gamma} _ {\chi} (s) f\right) (x) = \chi (s ^ {- 1}, x) f (s ^ {- 1} x).\tag{4}
\]

Then \(\gamma_{\chi}(e)f = f\) and

\[
\begin{array}{r l} & {\big (\pmb {\gamma} _ {\chi} (s) \pmb {\gamma} _ {\chi} (s ^ {\prime}) f \big) (x) = \chi (s ^ {- 1}, x) \big (\pmb {\gamma} _ {\chi} (s ^ {\prime}) f \big) (s ^ {- 1} x)} \\ & {\qquad = \chi (s ^ {- 1}, x) \chi (s ^ {- 1}, s ^ {- 1} x) f (s ^ {- 1} s ^ {- 1} x)} \\ & {\qquad = \chi \big ((s s ^ {\prime}) ^ {- 1}, x) \big) f \big ((s s ^ {\prime}) ^ {- 1} x \big) = \big (\pmb {\gamma} _ {\chi} (s s ^ {\prime}) f \big) (x),} \end{array}
\]

thus \(\pmb{\gamma}_{\chi}\) is a linear representation of G. For \(\chi = 1\) , one recovers the endomorphisms \(\pmb{\gamma}(s)\) (Ch. VII, §1, No.~1, formula (3)).

Suppose now that G and X are locally compact, G operating continuously on X, and \(\chi\) continuous on \(G \times X\) . Then \(\mathcal{C}(X)\) and \(\mathcal{K}(X)\) are stable for the \(\boldsymbol{\gamma}_{\chi}(s)\) , whence linear representations of G in \(\mathcal{C}(X)\) and \(\mathcal{K}(X)\) which we shall again denote \(\gamma_{\chi}\) .

PROPOSITION 4. --- The linear representations \(\pmb{\gamma}_{\chi}\) of G in \(\mathcal{C}(\mathbf{X})\) and \(\mathcal{K}(\mathbf{X})\) are continuous.

The mapping \((s, f) \mapsto (s, \gamma(s)f)\) of \(\mathbf{G} \times \mathcal{C}(\mathbf{X})\) into \(\mathbf{G} \times \mathcal{C}(\mathbf{X})\) is continuous (No.~2, Lemma 3). On the other hand, the mapping \((s, f) \mapsto \chi(s, \cdot)f\) of \(\mathbf{G} \times \mathcal{C}(\mathbf{X})\) into \(\mathcal{C}(\mathbf{X})\) is continuous; for, if \(s\) tends to \(s_0\) in \(\mathbf{G}\) , then \(\chi(s, \cdot)\) tends to \(\chi(s_0, \cdot)\) uniformly on every compact subset of \(\mathbf{X}\) ; if, moreover, \(f\) tends to \(f_0\) in \(\mathcal{C}(\mathbf{X})\) , then \(\chi(s, \cdot)f\) tends to \(\chi(s_0, \cdot)f_0\) uniformly on every compact subset of \(\mathbf{X}\) , whence our assertion. Thus the representation \(\gamma_{\chi}\) of \(\mathbf{G}\) in \(\mathcal{C}(\mathbf{X})\) is continuous.

Let us show that the representation \(\pmb{\gamma}_{\chi}\) of G in \(\mathcal{K}(\mathrm{X})\) is continuous. Since \(\mathcal{K}(\mathrm{X})\) is the direct limit of Banach spaces, it is barreled (TVS, III, §4, No.~1, Cor. 3 of Prop. 3), thus it suffices to prove that \(\pmb{\gamma}_{\chi}\) is separately continuous (No.~1, Prop. 1). Now, let H be a compact subset of X and let \(s_0 \in \mathrm{G}\) . Let V be a compact neighborhood of \(s_0\) in G, and let L = VH, which is compact in X. For every \(f \in \mathcal{K}(\mathrm{X}, \mathrm{H})\) , the support of \(\pmb{\gamma}_{\chi}(s_0)f\) is contained in L, and

\[
\sup _ {x \in X} | (\boldsymbol {\gamma} _ {\chi} (s _ {0}) f) (x) | \leqslant \sup _ {x \in L} | \chi (s _ {0} ^ {- 1}, x) | \cdot \sup _ {x \in X} | f (x) |,
\]

therefore \(f \mapsto \boldsymbol{\gamma}_{\chi}(s_0)f\) is a continuous linear mapping of \(\mathcal{K}(\mathrm{X},\mathrm{H})\) into \(\mathcal{K}(\mathrm{X},\mathrm{L})\) ; it follows that \(f \mapsto \boldsymbol{\gamma}_{\chi}(s_0)f\) is a continuous linear mapping of \(\mathcal{K}(\mathrm{X})\) into itself (TVS, II, §4, No.~4, Prop. 5). On the other hand, the topology of \(\mathcal{K}(\mathrm{X},\mathrm{L})\) is induced by that of \(\mathcal{C}(\mathrm{X})\) . By what has already been proved, the mapping \(s \mapsto \boldsymbol{\gamma}_{\chi}(s)f\) of \(\mathrm{V}\) into \(\mathcal{K}(\mathrm{X},\mathrm{L})\) is continuous. This completes the proof that \(\boldsymbol{\gamma}_{\chi}\) is separately continuous.

PROPOSITION 5. --- Suppose that each function \(\chi(s,\cdot)\) is bounded. Then \(\gamma_{\chi}\) leaves \(\overline{\mathcal{K}(X)}\) stable, and the linear representation \(\gamma_{\chi}\) of \(G\) in \(\overline{\mathcal{K}(X)}\) is continuous.

It is clear that \(\pmb{\gamma}_{\chi}\) leaves \(\overline{\mathcal{K}(\mathrm{X})}\) stable and that each of the \(\pmb{\gamma}_{\chi}(s)\) is continuous in \(\overline{\mathcal{K}(\mathrm{X})}\) . On the other hand, for every \(f \in \mathcal{K}(\mathrm{X})\) , \(s \mapsto \pmb{\gamma}_{\chi}(s)f\) is a continuous mapping of \(\mathbf{G}\) into \(\mathcal{K}(\mathbf{X})\) and a fortiori into \(\overline{\mathcal{K}(\mathbf{X})}\) . Therefore the representation \(\pmb{\gamma}_{\chi}\) in \(\overline{\mathcal{K}(\mathbf{X})}\) is continuous (No.~1, Prop. 2).

\subsection{Example: linear representations in spaces of measures}

Again let G be a locally compact group, operating continuously on the left in a locally compact space X, and let \(\chi\) be a continuous multiplier on \(\mathbf{G} \times \mathbf{X}\) . The linear representation \(\boldsymbol{\gamma}_{\chi}\) of \(\mathbf{G}\) in \(\mathcal{K}(\mathbf{X})\) admits a contragredient representation in \(\mathcal{M}(\mathbf{X})\) , which we shall again denote by \(\boldsymbol{\gamma}_{\chi}\) , and which is defined by the following formula (where \(\mu \in \mathcal{M}(\mathbf{X})\) , \(f \in \mathcal{K}(\mathbf{X})\) ):

\[
\langle \boldsymbol {\gamma} _ {\chi} (s) \mu , f \rangle = \langle \mu , \boldsymbol {\gamma} _ {\chi} (s ^ {- 1}) f \rangle = \langle \chi (s, \cdot) \cdot \mu , \boldsymbol {\gamma} (s ^ {- 1}) f \rangle = \langle \boldsymbol {\gamma} (s) (\chi (s, \cdot) \cdot \mu), f \rangle ,
\]

whence

\[
\boldsymbol {\gamma} _ {\chi} (s) \mu = \boldsymbol {\gamma} (s) (\chi (s, \cdot) \cdot \mu) = (\boldsymbol {\gamma} (s) \chi (s, \cdot)) \cdot (\boldsymbol {\gamma} (s) \mu).
\]

We note that

\[
\left(\gamma (s) \chi (s, \cdot)\right) (x) = \chi (s, s ^ {- 1} x).
\]

The linear representation \(\gamma_{\chi}\) of G in \(\mathcal{C}(X)\) admits a contragredient representation in the space \(\mathcal{C}'(X)\) of measures on X with compact support, a representation which we again denote by \(\gamma_{\chi}\) ; the endomorphisms \(\gamma_{\chi}(s)\) of \(\mathcal{C}'(X)\) are the restrictions of the endomorphisms \(\gamma_{\chi}(s)\) of \(\mathcal{M}(X)\) .

PROPOSITION 6. --- If one equips \(\mathcal{M}(\mathrm{X})\) (resp. \(\mathcal{C}'(\mathrm{X})\) ) with the topology of uniform convergence in the compact subsets of \(\mathcal{H}(\mathrm{X})\) (resp. \(\mathcal{C}(\mathrm{X})\) ), then the linear representation \(\gamma_{\chi}\) of \(\mathrm{G}\) in \(\mathcal{M}(\mathrm{X})\) (resp. \(\mathcal{C}'(\mathrm{X})\) ) is continuous.

PROPOSITION 7. --- Suppose that each function \(\chi(s,\cdot)\) is bounded. Then \(\gamma_{\chi}\) leaves stable \(\mathcal{M}^{1}(X)\) and, if \(\mathcal{M}^{1}(X)\) is equipped with the topology of uniform convergence in the compact subsets of \(\overline{\mathcal{K}(X)}\) , then the linear representation \(\gamma_{\chi}\) of \(G\) in \(\mathcal{M}^{1}(X)\) is continuous.

These propositions result from Props. 3, 4, 5.

\subsection{\texorpdfstring{Example: linear representations in the spaces \(L^{p}\)}{Example: linear representations in the spaces L\^{}\{p\}}}

Again let G be a locally compact group, operating continuously on the left in a locally compact space X. Let \(\beta\) be a positive measure on X with support X. Let us assume that there exists a continuous function \(\chi > 0\) on \(G \times X\) such that, for every \(s \in G\) ,

\[
\boldsymbol {\gamma} (s) \beta = \chi (s ^ {- 1}, \cdot) \cdot \beta
\]

(which implies in particular that \(\beta\) is quasi-invariant under \(\mathbf{G}\) ). Then, \(\chi\) is a multiplier. For, given \(s, t\) in \(\mathbf{G}\) , one has

\[
\begin{array}{r l} \boldsymbol {\gamma} (s) \boldsymbol {\gamma} (t) \beta = \boldsymbol {\gamma} (s) \big (\chi (t ^ {- 1}, \cdot) \cdot \beta \big) & = \big (\boldsymbol {\gamma} (s) \chi (t ^ {- 1}, \cdot) \big) \cdot \big (\boldsymbol {\gamma} (s) \beta \big) \\ & = \big (\boldsymbol {\gamma} (s) \chi (t ^ {- 1}, \cdot) \big) \cdot \chi (s ^ {- 1}, \cdot) \cdot \beta , \\ \boldsymbol {\gamma} (s t) \beta & = \chi (t ^ {- 1} s ^ {- 1}, \cdot) \cdot \beta , \end{array}
\]

therefore

\[
\chi (t ^ {- 1}, s ^ {- 1} x) \chi (s ^ {- 1}, x) = \chi (t ^ {- 1} s ^ {- 1}, x)
\]

locally \(\beta\) -almost everywhere, consequently everywhere, since \(\chi\) is continuous and \(\beta\) has support \(X\) .

Let \(p \in [1, +\infty[\) . For every \(f \in \mathcal{L}_{\mathbf{C}}^{p}(\mathbf{X}, \beta)\) and every \(s \in \mathbf{G}\) , let \(\gamma_{\chi,p}(s)f\) be the function on \(\mathbf{X}\) defined by

\[
\left(\gamma_ {\chi , p} (s) f\right) (x) = \chi (s ^ {- 1}, x) ^ {1 / p} f (s ^ {- 1} x).
\]

One has

\[
\begin{array}{c} \int^ {*} | \chi (s ^ {- 1}, x) ^ {1 / p} f (s ^ {- 1} x) | ^ {p} d \beta (x) = \int^ {*} | f (s ^ {- 1} x) | ^ {p} \chi (s ^ {- 1}, x) d \beta (x) \\ = \int | f (x) | ^ {p} d \beta (x), \end{array}
\]

therefore \(\gamma_{\chi,p}(s)f\in\mathcal{L}_{\mathbf{C}}^{p}(\mathrm{X},\beta)\) . One sees that \(\gamma_{\chi,p}(s)\) is an isometric endomorphism of \(\mathcal{L}_{\mathbf{C}}^{p}(\mathrm{X},\beta)\) and defines, by passage to the quotient, an isometric endomorphism of \(\mathrm{L}_{\mathbf{C}}^{p}(\mathrm{X},\beta)\) , also denoted \(\gamma_{\chi,p}(s)\) . On the other hand, \(\chi^{1/p}\) is obviously a multiplier, therefore \(\gamma_{\chi,p}\) is a linear representation of G in \(\mathrm{L}_{\mathbf{C}}^{p}(\mathrm{X},\beta)\) by what we have seen in No.~3.

PROPOSITION 8. --- The linear representation \(\gamma_{\chi,p}\) of G in \(\mathrm{L}_{\mathbf{C}}^{p}(\mathbf{X},\beta)\) is continuous and isometric.

Let \(f \in \mathcal{K}(\mathrm{X})\) . When \(s\) tends to \(s_0\) in \(\mathrm{G}\) , \(\boldsymbol{\gamma}_{\chi,p}(s)f\) tends to \(\boldsymbol{\gamma}_{\chi,p}(s_0)f\) in \(\mathcal{K}(\mathrm{X})\) , hence in \(\mathrm{L}_{\mathbf{C}}^p (\mathrm{X},\beta)\) . Since the \(\boldsymbol{\gamma}_{\chi,p}(s)\) are isometric, Prop. 8 is obtained by applying Remark 2 of No.~1.

For the case that \(\chi\) is not assumed continuous, cf.~§4, Exer. 13.

PROPOSITION 9. --- Suppose that each function \(\chi(s,\cdot)\) is bounded. Then \(\gamma_{\chi}\) leaves \(\mathrm{L}_{\mathbf{C}}^{p}(\mathrm{X},\beta)\) stable, and the linear representation \(\gamma_{\chi}\) of \(\mathbf{G}\) in \(\mathrm{L}_{\mathbf{C}}^{p}(\mathrm{X},\beta)\) is continuous.

Let \(f\in \mathcal{L}_{\mathbf{C}}^{p}(\mathrm{X},\beta)\) . Then

\[
\begin{array}{r l} \int^ {*} | \chi (s ^ {- 1}, x) f (s ^ {- 1} x) | ^ {p} d \beta (x) \\ & \leqslant \sup _ {x \in X} \chi (s ^ {- 1}, x) ^ {p - 1} \int^ {*} | f (s ^ {- 1} x) | ^ {p} \chi (s ^ {- 1}, x) d \beta (x) \\ & = \sup _ {x \in X} \chi (s ^ {- 1}, x) ^ {p - 1} \int | f (x) | ^ {p} d \beta (x), \end{array}
\]

therefore \(\pmb{\gamma}_{\chi}(s)f\in \mathcal{L}_{\mathbf{C}}^{p}(\mathrm{X},\beta)\) , and

\[
\left\| \boldsymbol {\gamma} _ {\chi} (s) \right\| \leqslant \sup _ {x \in X} \chi (s ^ {- 1}, x) ^ {1 / q},\tag{5}
\]

where \(q\) denotes the exponent conjugate to \(p\) . If \(f \in \mathcal{K}(\mathrm{X})\) , then \(\boldsymbol{\gamma}_{\chi}(s)f\) tends to \(\boldsymbol{\gamma}_{\chi}(s_0)f\) in \(\mathcal{K}(\mathrm{X})\) , hence in \(\mathcal{L}_{\mathbf{C}}^{p}(\mathrm{X},\beta)\) , as \(s\) tends to \(s_0\) . Therefore the representation \(\boldsymbol{\gamma}_{\chi}\) of \(\mathrm{G}\) in \(\mathrm{L}_{\mathbf{C}}^{p}(\mathrm{X},\beta)\) is continuous (No.~1, Prop. 2).

Properties analogous to those of Nos. 3, 4, 5 hold if G operates on the right in X.

In particular, if one regards G as operating on itself by left or right translations, and if one takes \(\chi = 1\) , one obtains the left and right regular representations of G in \(\mathcal{C}(\mathrm{G})\) , \(\mathcal{K}(\mathrm{G})\) , \(\overline{\mathcal{K}(\mathrm{G})}\) , \(\mathcal{C}'(\mathrm{G})\) , \(\mathcal{M}(\mathrm{G})\) , \(\mathcal{M}^{1}(\mathrm{G})\) . If one takes \(\beta\) to be a left (resp. right) Haar measure on G, and if one takes \(\chi = 1\) , one obtains the left (resp. right) regular representation of G in \(\mathrm{L}_{\mathbf{C}}^{p}(\mathrm{G}, \beta)\) .

\subsection{Extension of a linear representation of G to the measures on G}

Let G be a locally compact group, E a locally convex space, U a linear representation of G in E. Assume U to be continuous and E quasi-complete. Then, for every measure \(\mu \in \mathcal{C}'(G)\) , one has

\[
\int_ {\mathbf {G}} U (s) d \mu (s) \in \mathcal {L} (\mathbf {E}; \mathbf {E})
\]

(Ch. VI, §1, No.~7). We shall write \(U(\mu) = \int_{\mathbf{G}} U(s) d\mu(s)\) . We equip \(\mathcal{C}'(\mathbf{G})\) with the topology of compact convergence in \(\mathcal{C}(\mathbf{G})\) . The mapping \((\mu, x) \mapsto U(\mu)x\) of \(\mathcal{C}'(\mathbf{G}) \times \mathbf{E}\) into \(\mathbf{E}\) is hypocontinuous relative to the equicontinuous subsets of \(\mathcal{C}'(\mathbf{G})\) and the compact subsets of \(\mathbf{E}\) ; in particular, the mapping \(\mu \mapsto U(\mu)\) of \(\mathcal{C}'(\mathbf{G})\) into \(\mathcal{L}(\mathbf{E}; \mathbf{E})\) (equipped with the topology of compact convergence) is continuous (loc. cit., Prop. 16).

In order to be able to apply these results later on, we note that if X is a locally compact space then \(\mathcal{C}(\mathrm{X})\) , equipped with the topology of compact convergence, is complete (GT, X, §1, No.~6, Cor. 3 of Th. 2). On the other hand, \(\mathcal{K}(\mathrm{X})\) is barreled, therefore its dual \(\mathcal{M}(\mathrm{X})\) , equipped with the topology of compact convergence in \(\mathcal{K}(\mathrm{X})\) , is quasi-complete (TVS, III, §4, No.~2, Cor. 4 of Th. 1). Of course, \(\overline{\mathcal{K}(\mathrm{X})}\) is complete for the topology deduced from its norm, therefore its dual \(\mathcal{M}^{1}(\mathrm{X})\) is quasi-complete for the topology of compact convergence in \(\overline{\mathcal{K}(\mathrm{X})}\) (loc. cit.).

Let us now assume that U is a continuous linear representation of the locally compact group G on a Banach space E. Set \(g(s) = \|U(s)\|\) for all \(s \in G\) . Then, if \(\mu\) is a measure on G such that g is \(\mu\) -integrable, one has \(\int_{\mathbf{G}} U(s) d\mu(s) \in \mathcal{L}(\mathbf{E}; \mathbf{E})\) and \(\left\|\int_{\mathbf{G}} U(s) d\mu(s)\right\| \leqslant \int g(s) d|\mu|(s)\) (Ch. VI, §1, No.~7, Remark 1). We again write \(U(\mu) = \int_{\mathbf{G}} U(s) d\mu(s)\) .

\begin{enumerate}
\def\labelenumi{\arabic{enumi}.}
\setcounter{enumi}{6}
\tightlist
\item
  Relations between the endomorphisms \(U(\mu)\) and the endomorphisms \(U(s)\)
\end{enumerate}

Lemma 4. --- Let T be a locally compact space, a a point of T, M a subset of \(\mathcal{M}(\mathrm{T})\) , and \(\mathfrak{F}\) a filter on M. Assume that:

\begin{enumerate}
\def\labelenumi{(\roman{enumi})}
\item
  for every compact subset \(\mathbf{K}\) of \(\mathbf{T}\) , the numbers \(|\mu|(\mathbf{K})\) , for \(\mu \in \mathbf{M}\) , are bounded above;
\item
  \(\lim_{\mu, \mathfrak{F}} |\mu| (\mathrm{K}) = 0\) for every compact subset \(\mathrm{K}\) of \(\mathrm{T} - \{a\}\) .
\item
  there exists a compact neighborhood V of a in T such that \(\lim_{\mu, \mathfrak{F}} \mu(V) = 1\) .
\end{enumerate}

Then the filter \(\mathfrak{F}\) converges to \(\varepsilon_{a}\) in \(\mathcal{M}(\mathrm{T})\) equipped with the topology of compact convergence in \(\mathcal{H}(\mathrm{T})\) .

By the hypothesis (i), M is an equicontinuous subset of \(\mathcal{M}(\mathrm{T})\) since it is vaguely bounded and \(\mathcal{H}(\mathrm{T})\) is barreled (TVS, III, §4, No.~2, Th. 1). It therefore suffices (GT, X, §2, No.~4, Th. 1) to prove that if \(f \in \mathcal{H}(\mathrm{T})\) , then \(\lim_{\mu, \mathfrak{F}} \mu(f) = f(a)\) . Let K be the union of V and the support of \(f\) ; if \(\mathrm{K}'\) is the closure of \(\mathrm{K} - \mathrm{V}\) , one has

\[
| \mu (\mathrm{K}) - \mu (\mathrm{V}) | = | \mu (\mathrm{K} - \mathrm{V}) | \leqslant | \mu | (\mathrm{K} ^ {\prime});
\]

since \(\mathbf{K}'\) is compact and does not contain \(a\) , one concludes from this that \(\lim_{\mu, \mathfrak{F}} \mu(\mathbf{K}) = 1\) . Let \(\varepsilon > 0\) , and let \(\mathbf{W}\) be an open neighborhood of \(a\) in \(\mathbf{K}\) such that \(|f(t) - f(a)| \leqslant \varepsilon\) for \(t \in \mathbf{W}\) ; one can write

\[
\mu (f) - f (a) = f (a) \big (\mu (\mathrm{K}) - 1 \big) + \int_ {\mathrm{K}} \big (f (t) - f (a) \big) d \mu (t);
\]

the integral over K may be written as the sum of the analogous integrals over W and K - W; if \(C = \sup |f|\) , one therefore has

\[
| \mu (f) - f (a) | \leqslant \mathrm{C} | \mu (\mathrm{K}) - 1 | + \varepsilon \cdot | \mu | (\mathrm{K}) + 2 \mathrm{C} \cdot | \mu | (\mathrm{K} - \mathrm{W}).
\]

Since the first and third terms on the right side tend to 0 with respect to \(\mathfrak{F}\) , one sees that indeed \(\lim_{\mu \to \mathfrak{F}} \mu(f) = f(a)\) .

COROLLARY 1. --- With hypotheses as in Lemma 4, suppose in addition that there exists a compact subset \(K_{0}\) of T containing the supports of all of the measures \(\mu \in M\) . Then F also converges to \(\varepsilon_{a}\) in \(\mathcal{C}'(T)\) equipped with the topology of compact convergence in \(\mathcal{C}(T)\) .

For, the restriction mapping of \(\mathcal{C}(\mathrm{T})\) into \(\mathcal{C}(\mathrm{K}_{0})\) is continuous; therefore, if H is a compact subset of \(\mathcal{C}(\mathrm{T})\) , then the restrictions to \(K_{0}\) of the functions in H form a compact subset of \(\mathcal{C}(\mathrm{K}_{0})\) . It then suffices to apply Lemma 4 on replacing T by \(K_{0}\) .

COROLLARY 2. --- With hypotheses as in Cor. 1, let f be a continuous mapping of T into a quasi-complete locally convex space E. Then

\[
\lim _ {\mu , \mathfrak {F}} \int f (t) d \mu (t) = f (a).
\]

This follows from Cor. 1, and Prop. 14 of Ch. VI, §1, No.~6.

COROLLARY 3. --- Let G be a locally compact group, E a quasi-complete locally convex space, and U a continuous linear representation of G in E. Let \(\beta\) be a positive measure on G, a an element of G, and B a base for the filter of neighborhoods of a, formed of compact neighborhoods. For every \(V \in B\) , let \(f_{V}\) be a continuous function \(\geqslant 0\) on G, with support contained in V, and such that \(\int f_{V} d\beta = 1\) . Then, for every \(x \in E\) ,

\[
U (a) x = \lim _ {\mathrm{V}} U (f _ {\mathrm{V}} \cdot \beta) x,
\]

the limit being taken with respect to the section filter of \(\mathfrak{B}\) .

The mapping \(s \mapsto U(s)x\) of \(G\) into \(E\) is continuous. By Cor. 2, \(U(a)x = \lim_{V} \int (U(s)x) \cdot f_V(s) d\beta(s)\) with respect to the section filter of \(\mathfrak{B}\) , that is, \(U(a)x = \lim_{V} U(f_V \cdot \beta)x\) .

PROPOSITION 10. --- Let G be a locally compact group, E a quasi-complete locally convex space, U a continuous linear representation of G in E, and \(\beta\) a positive measure on G with support G.

\begin{enumerate}
\def\labelenumi{(\roman{enumi})}
\item
  The vectors \(U(f \cdot \beta)x\) , where \(f\) runs over \(\mathcal{K}(\mathbf{G})\) and \(x\) runs over \(\mathbf{E}\) , are dense in \(\mathbf{E}\) .
\item
  Let \(\mathbf{F}\) be a closed linear subspace of \(\mathbf{E}\) . If \(\mathbf{F}\) is stable for \(U\) , then \(U(\mu)(\mathbf{F}) \subset \mathbf{F}\) for every \(\mu \in \mathcal{C}'(\mathbf{G})\) . Conversely, if \(U(f \cdot \beta) \subset \mathbf{F}\) for every \(f \in \mathcal{K}(\mathbf{G})\) , then \(\mathbf{F}\) is stable for \(U\) .
\end{enumerate}

The first part of (ii) is immediate, since the restrictions of the \(U(s)\) to F ( \(s \in G\) ) define a continuous linear representation of G in the quasi-complete locally convex space F. The second part of (ii), and (i), follow from Cor. 3 of Lemma 4.

\section{CONVOLUTION OF MEASURES ON GROUPS}

\subsection{Algebras of measures}

Let G be a locally compact group. It will be understood, once and for all, that the measures \(\mu_{1},\ldots,\mu_{n}\) on G are said to be convolvable if they are so for the mapping

\[
(x _ {1}, x _ {2}, \ldots , x _ {n}) \mapsto x _ {1} x _ {2} \cdot \cdot \cdot x _ {n};
\]

and it is by means of this mapping that the convolution product \(\ast_{i}\mu_{i}\) will always be taken. If \(s\in G, t\in G\) , then

\[
\varepsilon_ {s} * \varepsilon_ {t} = \varepsilon_ {s t}.\tag{1}
\]

If \(s \in \mathbf{G}\) and \(\mu \in \mathcal{M}(\mathbf{G})\) , then

\begin{enumerate}
\def\labelenumi{(\arabic{enumi})}
\setcounter{enumi}{1}
\tightlist
\item
\end{enumerate}

\[
\begin{array}{l} \varepsilon_ {s} * \mu = \gamma (s) \mu \\ \mu * \varepsilon_ {s} = \delta (s ^ {- 1}) \mu \end{array}\tag{3}
\]

by §1, No.~1, Example 3. If G is abelian, to say that \(\mu_{1}\) and \(\mu_{2}\) are convolvable is equivalent to saying that \(\mu_{2}\) and \(\mu_{1}\) are convolvable, and one then has \(\mu_{1} * \mu_{2} = \mu_{2} * \mu_{1}\) . When G is not abelian, it can happen that \(\mu_{1}\) and \(\mu_{2}\) are convolvable, without \(\mu_{2}\) and \(\mu_{1}\) being so (Exer. 12).

PROPOSITION 1. --- Let G be a locally compact group, \(\lambda, \mu, \nu\) measures \(\neq 0\) on G.

\begin{enumerate}
\def\labelenumi{(\roman{enumi})}
\tightlist
\item
  If \(\lambda, \mu, \nu\) are convolvable, then so are \(\lambda\) and \(\mu\) , \(|\lambda| * |\mu|\) and \(\nu\) , \(\mu\) and \(\nu\) , \(\lambda\) and \(|\mu| * |\nu|\) , and one has
\end{enumerate}

\[
\lambda * \mu * \nu = (\lambda * \mu) * \nu = \lambda * (\mu * \nu).
\]

\begin{enumerate}
\def\labelenumi{(\roman{enumi})}
\setcounter{enumi}{1}
\tightlist
\item
  If \(\lambda\) and \(\mu\) are convolvable, as well as \(|\lambda| * |\mu|\) and \(\nu\) , then \(\lambda, \mu, \nu\) are convolvable. Similarly if \(\mu\) and \(\nu\) are convolvable, as well as \(\lambda\) and \(|\mu| * |\nu|\) .
\end{enumerate}

This follows from Prop. 1 of §1, No.~2.

\pandocbounded{\includegraphics[keepaspectratio]{images/part001_0fdec0aa.jpg}}

There can exist measures \(\lambda, \mu, \nu\) on \(G\) such that the convolution products \(\lambda * \mu\) , \((\lambda * \mu) * \nu\) , \(\mu * \nu\) , \(\lambda * (\mu * \nu)\) are all defined, and yet \((\lambda * \mu) * \nu \neq \lambda * (\mu * \nu)\) (cf.~Exer. 4).

Let \(\rho\) be a lower semi-continuous finite function \(>0\) on \(\mathbf{G}\) such that \(\rho(st) \leqslant \rho(s)\rho(t)\) for all \(s, t\) in \(\mathbf{G}\) . We denote by \(\mathcal{M}^{\rho}(\mathbf{G})\) the vector space of measures \(\lambda\) on \(\mathbf{G}\) such that \(\rho\) is \(\lambda\) -integrable, and by \(\|\lambda\|_{\rho}\) (or simply \(\|\lambda\|\) ) the norm \(\int_{\mathbf{G}} \rho(s)d|\lambda|(s)\) on this space. When \(\rho = 1\) , one recovers the set \(\mathcal{M}^{1}(\mathbf{G})\) of bounded measures on \(\mathbf{G}\) .

PROPOSITION 2. --- (i) Any two elements of \(\mathcal{M}^{\rho}(\mathrm{G})\) are convolvable.

\begin{enumerate}
\def\labelenumi{(\roman{enumi})}
\setcounter{enumi}{1}
\item
  For convolution, and for the norm \(\| \lambda \|\) , \(\mathcal{M}^{\rho}(\mathbf{G})\) is a complete normed algebra, admitting \(\varepsilon_{e}\) as unity element.
\item
  \(\mathcal{C}'(\mathbf{G})\) is a subalgebra of \(\mathcal{M}^{\rho}(\mathbf{G})\) .
\end{enumerate}

Let \(\lambda, \mu\) be in \(\mathcal{M}^{\rho}(\mathbf{G})\) , and let us show that \(\lambda\) and \(\mu\) are convolvable. Let \(f \in \mathcal{K}_{+}(\mathbf{G})\) . Since \(\rho\) is \(>0\) and lower semi-continuous, there exists a constant \(k > 0\) such that \(f \leqslant k\rho\) . Then

\[
\begin{array}{r l} & {\int^ {*} f (s t) d | \lambda | (s) d | \mu | (t) \leqslant k \int^ {*} \rho (s t) d | \lambda | (s) d | \mu | (t)} \\ & {\qquad \leqslant k \int^ {*} \rho (s) \rho (t) d | \lambda | (s) d | \mu | (t)} \\ & {\qquad = k \bigg (\int^ {*} \rho (s) d | \lambda | (s) \bigg) \bigg (\int^ {*} \rho (t) d | \mu | (t) \bigg)} \end{array}
\]

(Ch. V, §8, No.~3, Cor. 1 of Prop. 8). Therefore \((s,t) \mapsto f(st)\) is \((\lambda \otimes \mu)\) -integrable, so that \(\lambda\) and \(\mu\) are convolvable. On the other hand, using Ch. V (§1, Prop. 4, §6, Prop. 2, §8, Cor. 1 of Prop. 8) and the fact that \((s,t) \mapsto \rho(s)\rho(t)\) is lower semi-continuous in \(G \times G\) , one has

\[
\begin{array}{r l} & {\int_ {\mathrm{G}} ^ {*} \rho (s) d | \lambda * \mu | (s) = \int_ {\mathrm{G}} ^ {\bullet} \rho (s) d | \lambda * \mu | (s)} \\ & {\quad \leqslant \int_ {\mathrm{G} \times \mathrm{G}} ^ {\bullet} \rho (s t) d | \lambda | (s) d | \mu | (t) \leqslant \int_ {\mathrm{G} \times \mathrm{G}} ^ {\bullet} \rho (s) \rho (t) d | \lambda | (s) d | \mu | (t)} \\ & {\quad = \int_ {\mathrm{G} \times \mathrm{G}} ^ {*} \rho (s) \rho (t) d | \lambda | (s) d | \mu | (t) = \| \lambda \| \cdot \| \mu \|.} \end{array}
\]

One sees that \(\lambda * \mu \in \mathcal{M}^{\rho}(\mathrm{G})\) and that \(\| \lambda * \mu \| \leqslant \| \lambda \| \cdot \| \mu \|\) . In view of Prop. 1, \(\mathcal{M}^{\rho}(\mathrm{G})\) is an algebra. The mapping \(\lambda \mapsto \rho \cdot \lambda\) is an isometric linear mapping \(\theta\) of \(\mathcal{M}^{\rho}(\mathrm{G})\) into \(\mathcal{M}^{1}(\mathrm{G})\) ; if \(\mu \in \mathcal{M}^{1}(\mathrm{G})\) then \(1 / \rho\) , which is locally bounded and upper semi-continuous, is locally \(\mu\) -integrable, and \(\rho\) is \((1 / \rho) \cdot \mu\) -integrable, thus \((1 / \rho) \cdot \mu \in \mathcal{M}^{\rho}(\mathrm{G})\) ; this proves that \(\theta\) is surjective; therefore \(\mathcal{M}^{\rho}(\mathrm{G})\) is a complete normed algebra. Finally, it is clear that \(\varepsilon_{e}\) is a unity element for \(\mathcal{M}^{\rho}(\mathrm{G})\) and that \(\mathcal{C}'(\mathrm{G})\) is a subalgebra of \(\mathcal{M}^{\rho}(\mathrm{G})\) (§1, No.~4, Cor. of Prop. 5).

If \(\rho = 1\) , Prop. 2, (i) and (ii) also follow from §1, Prop. 2.

PROPOSITION 3. --- Let \(\mu_1, \ldots, \mu_n\) be measures on G. If all of the \(\mu_i\) , except at most one, have compact support, then the \(\mu_i\) are convolvable.

For, let \(S_{i}\) be the support of \(\mu_{i}\) , and suppose that \(S_{i}\) is compact for \(i \neq i_{0}\) . Let K be a compact subset of G. The set of \((x_{1}, \ldots, x_{n}) \in \prod_{i} S_{i}\) such that \(x_{1}x_{2} \cdots x_{n} \in K\) is compact, because the conditions \(x_{i} \in S_{i}\) for all \(i, x_{1}x_{2} \cdots x_{n} \in K\) imply

\[
x _ {i _ {0}} \in \mathrm{S} _ {i _ {0} - 1} ^ {- 1} \dots \mathrm{S} _ {1} ^ {- 1} \mathrm{KS} _ {n} ^ {- 1} \dots \mathrm{S} _ {i _ {0} + 1} ^ {- 1}.
\]

Therefore the \(\mu_{i}\) are convolvable (§1, No.~4, Prop. 4).

PROPOSITION 4. --- The mapping \((\lambda, \mu) \mapsto \lambda * \mu\) (resp. \((\lambda, \mu) \mapsto \mu * \lambda)\) , where \(\lambda \in \mathcal{C}'(G)\) , \(\mu \in \mathcal{M}(G)\) , defines on \(\mathcal{M}(G)\) the structure of a left (resp. right) module over the algebra \(\mathcal{C}'(G)\) .

This follows from Props. 1 and 3.

PROPOSITION 5. --- Let \(\lambda\) be a left (resp. right) Haar measure on G, and \(\mu \in \mathcal{M}^1(\mathrm{G})\) . Then \(\mu\) and \(\lambda\) (resp. \(\lambda\) and \(\mu\) ) are convolvable, and \(\mu * \lambda = \mu(1)\lambda\) (resp. \(\lambda * \mu = \mu(1)\lambda\) ).

We can suppose that \(\mu \geqslant 0\) . Let \(f \in \mathcal{K}_{+}(\mathbf{G})\) . When \(\lambda\) is a left Haar measure,

\[
\int^ {*} d \mu (x) \int^ {*} f (x y) d \lambda (y) = \int^ {*} d \mu (x) \int f (y) d \lambda (y) = \lambda (f) \| \mu \|,
\]

therefore the function \((x,y)\mapsto f(xy)\) is \((\mu \otimes \lambda)\) -integrable, and its integral for \(\mu \otimes \lambda\) is \(\lambda (f)\| \mu \|\) . One argues similarly when \(\lambda\) is a right Haar measure.

PROPOSITION 6. --- Let \(\mu\) and \(\nu\) be two convolvable measures on \(\mathbf{G}\) . Let \(\chi\) be a continuous representation of \(\mathbf{G}\) in \(\mathbf{C}^*\) . Then \(\chi \cdot \mu\) and \(\chi \cdot \nu\) are convolvable and \((\chi \cdot \mu) * (\chi \cdot \nu) = \chi \cdot (\mu * \nu)\) .

Let \(f \in \mathcal{K}(\mathbf{G})\) . Then \(f\chi \in \mathcal{K}(\mathbf{G})\) , therefore the function

\[
(x, y) \mapsto f (x y) \chi (x y) = f (x y) \chi (x) \chi (y)
\]

on \(\mathbf{G} \times \mathbf{G}\) is integrable for \(\mu \otimes \nu\) ; therefore the function \((x, y) \mapsto f(xy)\) is integrable for \((\chi \cdot \mu) \otimes (\chi \cdot \nu)\) ; therefore \(\chi \cdot \mu\) and \(\chi \cdot \nu\) are convolvable. Moreover,

\[
\begin{array}{r l} \langle \chi \cdot \mu * \chi \cdot \nu , f \rangle & = \int f (x y) \chi (x) \chi (y) d \mu (x) d \nu (y) \\ & = \int (f \chi) (x y) d \mu (x) d \nu (y) = \langle \mu * \nu , \chi f \rangle , \end{array}
\]

whence \((\chi \cdot \mu) * (\chi \cdot \nu) = \chi \cdot (\mu * \nu)\) .

PROPOSITION 7. --- Let G and G' be two locally compact groups, u a continuous representation of G in G', and \(\mu_{1},\ldots,\mu_{n}\) measures on G, all \(\neq0\) . Then the following assertions are equivalent:

\begin{enumerate}
\def\labelenumi{(\roman{enumi})}
\item
  \(u\) is \(\mu_{i}\) -proper for all \(i\) , and the measures \(u(|\mu_i|)\) are convolvable;
\item
  the \(\mu_{i}\) are convolvable and \(u\) is proper for \(*_{i}(|\mu_{i}|)\)
\end{enumerate}

When these conditions are satisfied,

\[
\underset {i} {\ast} u (\mu_ {i}) = u \Big (\underset {i} {\ast} \mu_ {i} \Big).
\]

This follows from §1, No.~2, Cor. of Prop. 1.

COROLLARY. --- Let G be a locally compact group, \(\mu_{1},\ldots,\mu_{n}\) measures on G. For the sequence \((\mu_{i})_{1\leqslant i\leqslant n}\) to be convolvable, it is necessary and sufficient that the sequence \((\breve{\mu}_{n-i})_{0\leqslant i\leqslant n-1}\) be so, in which case

\[
\left(\mu_ {1} * \dots * \mu_ {n}\right) ^ {\vee} = \stackrel {\vee} {\mu} _ {n} * \dots * \stackrel {\vee} {\mu} _ {1}.
\]

This follows from Prop. 7 on considering the isomorphism \(x \mapsto x^{-1}\) of \(G\) onto the opposite group.

\subsection{The case of a group operating on a space}

Let X be a locally compact space on which a locally compact group G operates on the left continuously by

\[
(s, x) \mapsto s \cdot x.
\]

If \(\mu_{1},\ldots,\mu_{n}\) are measures on G and \(\nu\) is a measure on X, these will be said to be convolvable if they are so for the mapping \((s_{1},\ldots,s_{n},x)\mapsto s_{1}\cdots s_{n}x\) of \(G^{n}\times X\) into X, and their convolution product is to be understood in the sense of this mapping.

If \(s\in \mathbf{G}\) and \(x\in \mathbf{X}\) , then

\[
\varepsilon_ {s} * \varepsilon_ {x} = \varepsilon_ {s x}.\tag{4}
\]

If \(s \in \mathbf{G}\) and \(\mu \in \mathcal{M}(\mathbf{X})\) , then

\[
\varepsilon_ {s} * \mu = \gamma (s) \mu\tag{5}
\]

by §1, No.~1, Example 3.

PROPOSITION 8. --- Let \(\mu\) be a measure on G, \(\nu\) a measure on X.

\begin{enumerate}
\def\labelenumi{(\roman{enumi})}
\item
  If \(\mu\) has compact support, then \(\mu\) and \(\nu\) are convolvable.
\item
  If \(\nu\) has compact support, and if \(G\) operates properly in \(X\) , then \(\mu\) and \(\nu\) are convolvable.
\end{enumerate}

This follows from Prop. 4 of §1, No.~4.

PROPOSITION 9. --- For convolution, \(\mathcal{M}^{1}(\mathrm{X})\) is a left module over \(\mathcal{M}^{1}(\mathrm{G})\) , while \(\mathcal{M}(\mathrm{X})\) and \(\mathcal{C}'(\mathrm{X})\) are left modules over \(\mathcal{C}'(\mathrm{G})\) .

This follows from Prop. 8, and from §1, Props. 1, 3 and the Cor. of Prop. 5.

PROPOSITION 10. --- Let \(\mu\) be a measure on \(G\) , \(\nu\) a measure on \(X\) , \(\mu\) and \(\nu\) being convolvable. Suppose in addition that there exists a positive measure \(\beta\) on \(X\) such that \(\gamma(s)\nu\) has base \(\beta\) for every \(s \in G\) . Then \(\mu * \nu\) has base \(\beta\) .

Let K be a \(\beta\) -negligible compact subset of X. Then K is \(\pmb{\gamma}(s)|\nu|\) -negligible for every \(s \in G\) . Now,

\[
| \mu | * | \nu | = \int_ {\mathrm{G}} (\varepsilon_ {s} * | \nu |) d | \mu | (s)
\]

(§1, No.~5, Prop. 7), and the mapping \(s \mapsto \varepsilon_s * |\nu|\) is vaguely continuous (§2, Prop. 6). Therefore K is \(|\mu| * |\nu|\) -negligible by Ch. V, §3, No.~3, Th. 1. Therefore \(|\mu| * |\nu|\) has base \(\beta\) (Ch. V, §5, No.~5, Th. 2).
